\documentclass[../../book-main.tex]{subfiles}
\usepackage[UTF8]{ctex}
\graphicspath{{\subfix{../..}}}

\begin{document}

\chapter{引言}
\label{ch:intro}

\begin{quote}
\hfill    ``{\em 数学是给不同的事物赋予相同名称的艺术}.''

$~$ \hfill --- 埃内斯特·波利亚
\end{quote}
\vspace{5mm}

\section{智能、控制论、人工智能}
我们所处的世界既不是完全随机的，也不是完全不可预测的。\footnote{请注意，如果世界是完全随机的，那么智能生物就无需学习或记忆任何东西。}相反，它遵循一定的秩序、模式和规律，使其在很大程度上可以预测。\footnote{有些是确定性的，有些是概率性的。}生命的出现和存在依赖于一个可预测的生活环境。只有通过学习和记忆环境中的可预测部分，生命才能生存和繁衍，因为良好的决策和行动依赖于可靠的预测。由于世界中似乎有无限多的可预测事物，智能生物，如动物和人类，通过进化不断改进其探索和利用这种可预测性的能力，以获得更好的生活。为此，它们发展了越来越敏锐的感觉，包括视觉、听觉、触觉、味觉和嗅觉，以从这些高通量的感觉数据中感知外部环境中的可预测部分。因此，所有智能生物的基本任务之一是能够：
\begin{center}
    {\em 从大量感知数据中学习和记忆可预测的信息。}
\end{center}
在我们开始理解这是如何实现之前，我们需要考虑一些相关的问题：
\begin{itemize}
    \item 如何在数据中建模和表示这种可预测的信息？
    \item 如何有效地从数据中学习这种信息？
    \item 如何最好地组织或记忆这种信息以供将来使用？
\end{itemize}
本书旨在提供这些问题的一些答案。这些答案将帮助我们更好地理解智能，至少是使智能成为可能的计算原理和机制。有理由相信，现代科学作为最高形式的智能，遵循相同的原理和机制来进化。

\paragraph{自然界中智能的进化。}
%\sdb{对于本书的第一部分，这种组织方式与之前的组织方式（从可预测性开始）相比显得有些枯燥。从可预测性开始，我们从三个深刻而重要的问题开始（这些问题是非主流的），简洁地陈述，然后立即从基本原则开始理解它们；而这里我们开始得有些“教条”。这可能只是写作方式的问题，而不是组织结构的问题。}
%\yima{同意。我计划将此处的介绍移回第1.2节，作为本书的介绍。}

大约37亿年前地球上生命出现的一个必要条件是地球环境在很大程度上是可预测的。在环境中，生命发展了机制，使其能够学习可预测的部分并利用这些信息来生存。我们将这种能力称为{\em 智能}。在很大程度上，生命的进化也是智能的进化。在自然界中，智能主要通过两种类型的学习机制发展：{\em 系统发生}和{\em 个体发生} \cite{Wiener-Cybernetics-1961}。%\footnote{实际上，人类有一种独特的第三种智能，它被认为与这两种智能有很大不同。人类智能正式称为“人工智能”（AI）。我们将在后面讨论它。这里我们仅限于讨论自然界中普遍存在的智能，包括植物和动物。}

{\em 亲缘智能}是指通过物种的进化来学习。物种主要依靠从父母基因中继承的知识来继承和生存\footnote{自然的预训练模型}。物种通过基因的（随机）突变进化，然后被自然选择——适者生存。亲缘学习的主要特点是，个体的学习能力有限。学习是在物种层面通过自然选择的试错机制进行的，这可能被称为自然的强化学习。然而，这种强化学习过程可能极其缓慢、代价高昂且不可预测\footnote{已知，大约35亿年前，第一个生命形式——卢卡出现，生命依靠这种形式的进化大约30亿年，直到大约5.5亿年前个体神经系统出现}。

\begin{figure}
    \centering
\includegraphics[width=0.5\linewidth]{figs_chap1/DNAs.png}
%\includegraphics[width=0.28\linewidth]{figs_chap1/Luca.jpg}
\includegraphics[width=0.40\linewidth]{figs_chap1/Evolution.jpg}
    \caption{亲缘智能的进化：外部世界的知识通过DNA编码和传递（左），然后从DNA解码为RNA和蛋白质等。在生命进化的早期阶段（右），智能通过基因的（随机）突变和自然选择在物种层面发展知识——“适者生存”，这可以被视为一种原始形式的强化学习。}
    \label{fig:phylogenetic}
\end{figure}

{\em 个体智能}是指个体通过其感官信息和行为在特定生活环境中的学习机制，以及改进和适应其行为。个体智能的学习机制在大约5.5亿至6亿年前（在类似蠕虫的生物中）神经系统出现后成为可能。也就是说，有了神经系统和感官\footnote{大约与神经系统同时出现}，个体除了从DNA或基因中继承的知识外，还能够形成自己对世界的知识，即记忆。这种能力显著提高了个体的生存能力，并导致了大约5.3亿年前寒武纪时期的生命形式爆炸。与亲缘学习相比，个体学习速度更快、更经济且可控\footnote{可以在个体的时间和资源范围内实现}。

\begin{figure}
    \centering
\includegraphics[height=0.19\linewidth]{figs_chap1/Luca.jpg}
\includegraphics[height=0.19\linewidth]{figs_chap1/Worm.jpg}
\includegraphics[height=0.19\linewidth]{figs_chap1/Cambrian.jpg}
    \caption{生命进化，从大约35亿年前的第一个生命形式——类似单细胞的生物卢卡，到大约5.5亿年前类似蠕虫物种中第一个神经系统的出现，再到大约5.3亿年前寒武纪时期的生命形式爆炸。}
    \label{fig:evolution}
\end{figure}

需要注意的是，这两种学习机制都依赖于反馈（来自外部环境），即物种或个体行为的惩罚（死亡）或奖励（食物）来学习。因此，所有智能生命，无论是物种还是个体，都依赖于闭环反馈机制来学习和改进。我们还注意到，从植物到鱼类、鸟类再到哺乳动物，更高级的物种越来越依赖于其个体学习能力。它们出生后与父母在一起的时间越来越长，因为同一物种的个体需要在非常多样化的环境中生存。

\paragraph{人类智力的演变。}
自从大约31.5万年前智人出现以来，一种新的、更高层次的智力形式开始出现，并以更快的速度进化。几千年前，语言首先被口头使用\footnote{人们认为梵语是最早的语言之一，可以追溯到公元前5000年。}，然后被书写\footnote{苏美尔语被认为是现存最古老的文字之一，最早在公元前3100年左右的美索不达米亚南部被发现。}。这使得个体人类能够与他人交流和分享他们所学到的有用信息。因此，人类社区或社会可以像一个单一的智能有机体一样运作，学习速度更快，知识量更多，超过单个大脑所能容纳的。在某种程度上，书面语言或文本的作用类似于DNA和基因，它们允许人类社会积累并传递世界知识给下一代。科学可以被视为这种知识的一部分。我们可以将这种智力类型称为{\em 社会智力}，以区别于物种的系统发育智力和个体的个体发育智力。

\begin{figure}
    \centering
    \includegraphics[height=0.3\linewidth]{figs_chap1/Spoken-language.jpg}
   \hspace{3mm} \includegraphics[height=0.3\linewidth]{figs_chap1/Cuneiform.png}
   \hspace{3mm} \includegraphics[height=0.32\linewidth]{figs_chap1/Aristotle.jpg}
    \caption{口头交流和语言的发展（公元前10000-5000年）、书面语言（约公元前3000年）和数学（约公元前500年）的发展标志着人类智力演化的三个关键里程碑。}
    \label{fig:human-intelligence}
\end{figure}

大约两千年前，人类智力又经历了一次奇迹般的飞跃，使得哲学家和数学家能够发展出与经验观察和经验不直接相关的知识，例如数学抽象\footnote{数字、空间和时间的概念。}和逻辑推理。发展这些高级形式知识的能力被认为是人类独有的，并被称为“人工智能”（AI），这一术语由约翰·麦卡锡在1956年的达特茅斯暑期研讨会上提出。

因此，从我们所了解的自然中，从现在开始，每当使用“智力”这个词时，我们需要非常明确地指出我们指的是哪个层次的智力：

\begin{equation}
\mbox{\textbf{系统发育}} \;
   \Longrightarrow \; \mbox{\textbf{个体发育}} \; \Longrightarrow \; 
   \mbox{\textbf{社会}} \; \Longrightarrow \; 
   \mbox{\textbf{人工}}.
\end{equation}

明确的特征和区分是必要的和重要的，因为如果我们想将智力作为科学或数学主题进行研究，很可能每个层次的智力所涉及的学习和计算机制都是不同的，尽管无疑存在联系。

%在生物学领域，洛伦兹和廷伯根是研究动物行为的先驱（尽管不是智力本身）。在认知科学中，有一个悠久的历史，包括行为主义者（强化学习）、格式塔心理学家等。

%科学的启发。诺伯特·维纳的传记：“信息时代的黑暗英雄：寻找诺伯特·维纳——控制论之父”由康威和西格尔曼撰写。与许多生物学家的联系。维纳在学生时代学习过动物学。

\paragraph{机器智能的起源。}
在20世纪40年代，部分由于战争的努力，自然界中的智能启发了20世纪40年代的科学家们通过机器来模仿动物智能，这导致了诺伯特·维纳倡导的“控制论”运动。维纳在哈佛大学学习动物学作为本科生，但后来成为了一名数学家和控制理论家。维纳一生致力于理解和开发能够模仿动物智能行为的自主系统。今天，控制论计划通常被人们狭义地解释为主要关于反馈控制系统，维纳确实为此做出了最重要的技术贡献。但控制论计划比这更广泛和深刻。它更多地关注对智能的整体理解\footnote{至少在动物的水平上。}，并且实际上影响了一整代著名科学家的工作，包括沃伦·麦卡洛克、沃尔特·皮茨、克劳德·香农、约翰·冯·诺伊曼和艾伦·图灵。

\begin{figure}
    \centering
    \includegraphics[height=0.4\linewidth]{figures/Cybernetics1.jpg}
    \hspace{10mm} \includegraphics[height=0.4\linewidth]{figures/Cybernetics2.jpg}
    \caption{诺伯特·维纳于1948年出版的《控制论》一书\cite{Wiener-Cybernetics-1948}（左）和1961年的第二版\cite{Wiener-Cybernetics-1961}（右）。}
    \label{fig:cybernetcis}
\end{figure}

维纳可以说是第一个将智能视为{\em 系统}的人，而不是仅仅关注其某个组成部分或方面。他对智能的全面观点在1948年出版的著名著作《控制论：或动物和机器中的控制和通信》\cite{Wiener-Cybernetics-1948}中得到了详细阐述。在这本书及其1961年的第二版\cite{Wiener-Cybernetics-1961}（见图\ref{fig:cybernetcis}）中，他试图识别智能系统的一些必要特征和机制，包括但不限于：
\begin{itemize}
    \item 如何{\em 测量和存储}信息（在大脑中）以及如何与他人交流。 \footnote{诺伯特·维纳是第一个指出“信息”不是物质也不是能量，而是独立的研究量的人。} 这导致了克劳德·香农在1948年提出了信息论和编码理论。
    \item 如何根据现有信息{\em 纠正错误} 在预测和估计中。诺伯特·维纳本人在20世纪40年代帮助形式化了基于闭环反馈的控制系统理论。
    \item 如何从与潜在不合作对手或对抗环境的互动中学习{\em 做出更好的决策}。这在1944年由约翰·冯·诺伊曼形式化为博弈论。
\end{itemize}
1943年，受到维纳的控制论计划的极大启发，精神病学家沃伦·麦卡洛克和逻辑学家沃尔特·皮茨一起形式化了第一个神经元的计算模型\cite{McCulloch-Pitts}，称为{\em 人工神经元}，如图\ref{fig:neuron}所示。基于这个模型，在20世纪50年代，弗兰克·罗森布拉特建造了一个物理机器，名为{\em Mark I 感知器}，由数百个人工神经元组成的网络。感知器是第一个物理实现的人工神经网络，见图\ref{fig:perceptron}。值得注意的是，约翰·冯·诺伊曼在1945年提出的通用计算机架构也是为了促进构建能够物理实现控制论计划所建议机制的{\em 计算机器}的目标。

敏锐的读者可能已经注意到，20世纪40年代是一个真正的神奇时代：那个时代诞生了许多基本思想和有影响力的理论，包括神经元的数学模型、人工神经网络、信息论、控制论、博弈论和计算机器。我们现在知道，上述每项工作都成为了接下来几十年中某个科学或工程领域的基础。所有这些基本发现或理论都是受到尝试开发能够模仿自然界智能的机器这一目标的启发和推动。根据历史记录，维纳的控制论运动几乎影响了所有这些人和工作。

\begin{figure}
    \centering
    \includegraphics[height=0.3\linewidth]{figs_chap1/Wiener.png}
    \includegraphics[height=0.3\linewidth]{figs_chap1/Shannon.jpg}
    \includegraphics[height=0.3\linewidth]{figs_chap1/VonNeumann.jpg}
    \includegraphics[height=0.3\linewidth]{figs_chap1/Turing.jpeg}
    \caption{诺伯特·维纳（控制论和控制理论）、克劳德·香农（信息论）、约翰·冯·诺伊曼（博弈论）和艾伦·图灵（计算理论）。}
    \label{fig:god-fathers}
\end{figure}

尽管维纳在他的工作中已经识别出许多智能的关键特性和机制，但没有迹象表明他知道如何正确地将所有这些机制整合在一起，以构建一个完整的自主智能系统。根据今天的知识判断，他对这些机制的一些观点并不完全准确或完整。特别是，在《控制论》第二版的最后一章中 \cite{Wiener-Cybernetics-1961}，他指出，如果机器学习系统旨在模仿自然界中的典型学习机制，那么{\em 处理非线性}至关重要。但他没有提供任何具体和有效的解决方案来解决这一难题。为了辩护，当时很少有人知道如何处理，因为即使处理线性模型和系统的理论也还处于起步阶段。

尽管如此，我们还是不得不惊叹于维纳对非线性重要性的远见。正如我们将在本书中看到的，答案只是最近才被找到：通过深度神经网络实现的逐步线性和转换，可以有效地处理非线性（见第 \ref{ch:representation} 章）。此外，我们将在本书中尝试展示如何将上述所有机制自然地整合到一个完整的系统中，该系统将表现出自主智能系统的特征（见第 \ref{ch:autoencoding} 和 \ref{ch:closed-loop} 章）。

\paragraph{人工智能的起源。}
从维纳的《控制论》一书的副标题：{\em ``动物和机器中的控制和通信''}，可以看出20世纪40年代的研究主要旨在模仿动物层面的智能。正如我们之前提到的，20世纪40年代关于智能的研究议程在很大程度上受到诺伯特·维纳的控制论运动的影响。

艾伦·图灵是最早注意到这一局限性的人之一。在他的著名论文《{\em 计算机器与智能}》 \cite{Turing-1950} 中，图灵正式提出了机器是否能够模仿甚至达到人类水平的智能的问题，以至于无法与人类的智能能力区分开来。这现在被称为{\em 图灵测试}。

大约在1955年，一群年轻而雄心勃勃的科学家试图摆脱当时主导的控制论运动和研究议程，以便有机会创造自己的遗产。他们决定接受图灵模仿人类智能的挑战，并提议在1956年夏季于达特茅斯学院举办一个研讨会。他们在提案中明确表达了他们的意图：

\begin{quote}
    ``{\em 这项研究将基于这样一个假设进行：学习的各个方面或智能的任何其他特征原则上都可以被如此精确地描述，以至于可以制造出一台机器来模拟它。我们将尝试找出如何使机器使用语言、形成抽象和概念、解决目前保留给人类的问题类型，并自我改进}。''
\end{quote}

本质上，他们希望形式化并研究区分人类与动物的高级智能。他们考虑的主题包括抽象、符号方法、自然语言、演绎方法（包括因果推理和逻辑演绎等）。研讨会的组织者约翰·麦卡锡，当时是达特茅斯学院数学系的年轻助理教授，创造了现在著名的术语“人工智能”（AI），以概括那些被认为{\em 独特于人类智能}的特征或机制。

\paragraph{“人工智能”或“控制论”的复兴？}
读者可能已经知道，在过去十年左右，智能技术经历了爆炸式的发展，主要由深度人工神经网络的实践推动，这一发展始于2012年杰弗里·辛顿及其学生的工作\cite{Krizhevsky2012imagenet}。这个时代也被誉为“人工智能”（AI）的“复兴”。然而，从人们实际尝试解决的任务（识别、生成和预测）和技术上，我们实际上只是在模仿生命和动物智能中普遍存在的机制。即使在这一点上，严格来说，当前的“AI”模型和系统尚未完全实现20世纪40年代控制论运动中已知的所有必要机制。

因此，从技术上讲，过去十年在“人工智能”名义下取得的进展在很大程度上并没有服务于1956年达特茅斯研讨会所提出的“人工智能”计划的主要目的。相反，迄今为止主要实现的是与诺伯特·维纳在20世纪40年代提出的经典控制论计划的目标更为接近。将当前时代称为“控制论的复兴”可能更为恰当。\footnote{最近所谓的“具身AI”在自主智能机器人中的兴起与控制论计划的目标有更多相似之处。}只有当我们从科学和数学的角度充分理解我们真正做了什么，我们才能知道还有哪些工作需要完成，以及应该朝着哪个方向追求智能的本质。这是本书的主要目的之一。

\section{学习什么？}
\label{sec:what-to-learn}
%Notes: 智能的概念可以非常广泛和宽泛。但这也往往使其变得模糊不清。在这里，我们尝试给出一个更精确的描述，说明它主要是什么。我们的定义可能显著缩小了范围，但使其更易于实现、验证，甚至可测量（通过可计算的方式）。

\subsection{可预测性}
\label{sec:predictability}
携带可预测信息的数据以多种形式出现。在最自然的形式中，它们可以被建模为可预测和可计算的序列。可预测和可计算序列的概念及其性质是计算理论的核心，并且极大地推动了计算机的发明 \cite{Turing-1936}。可预测序列在（归纳）推理中的作用在20世纪60年代由Ray Solomonoff、Andrey Kolmogorov等人研究，作为对Claude Shannon的经典信息理论的推广 \cite{Kolmogorov1998OnTO}。为了理解可预测序列的概念，让我们首先从一些具体的简单例子开始。

\paragraph{标量情况。} 最简单的可预测离散序列可能是自然数序列：
\begin{equation}
   {S} =  1, 2, 3, 4, 5, 6, \ldots, n, n+1, \ldots
\end{equation}
其中下一个数字 $x_{n+1}$ 被定义为前一个数字 $x_n$ 加上1：
\begin{equation}
x_{n+1} = x_n + 1.    
\end{equation}
可以将可预测性的概念推广到任何序列 $\{x_n\}_{n=1}^\infty$，其中 $ x_n \in \mathbb{R}$，如果下一个数字 $x_{n+1}$ 总是可以从前一个数字 $x_n$ 计算出来：
\begin{equation}
    x_{n+1} = f(x_{n}), \quad x_n \in \mathbb{R}, \; n =  1, 2, 3, \ldots
\end{equation}
其中 $f(\cdot)$ 是一个可计算的（标量）函数。\footnote{这里我们强调函数 $f(\cdot)$ 本身是可计算的，例如，它可以在计算机上实现为一个程序。} 注意，这里我们强调函数 $f(\cdot)$ 必须是可计算的。有许多可以定义但不可计算的函数。Alan Turing在1936年的开创性工作 \cite{Turing-1936} 给出了可计算性的严格定义。在实践中，我们通常进一步假设 $f$ 是高效可计算的，并且具有良好的性质，例如连续性和可微性等。一旦我们对更精细的可计算性概念有了更多的了解，以及它们在机器学习和智能中的作用，这些性质的必要性将变得清晰。

\paragraph{多变量情况。}
当然，下一个数字的值也可以依赖于其前两个数字。例如，著名的斐波那契数列定义为：
\begin{equation}
    {S} = 1, 1, 2, 3, 5, 8, 13, 21, 34, 55, \ldots
\end{equation}
其中很容易看出：
\begin{equation}
    x_{n+2} = x_{n+1} + x_{n}, \quad  x_n \in \mathbb{R}, \;  n = 1, 2, 3, \ldots
\end{equation}
类似地，我们可以将这种递归推广到：
\begin{equation}
    x_{n+2} = f(x_{n+1}, x_{n}), \quad x_n \in \mathbb{R}, \;  n =  1, 2, 3, \ldots
\end{equation}
其中 $f(\cdot,\cdot)$ 是任何接受两个变量作为输入的可计算函数。我们可以进一步将可预测性的概念推广到下一个值依赖于其前 $d$ 个数字的序列：
\begin{equation}
    x_{n+d} = f(x_{n+d-1}, \ldots,  x_{n}), \quad  x_n \in \mathbb{R}, \; n =  1, 2, 3, \ldots
    \label{eqn:recursive-d}
\end{equation}
用于预测的前驱数字的数量 $d$ 称为递归预测的“阶数”。上述表达式 \eqref{eqn:recursive-d} 也称为“（自）回归”。这样的序列也称为“自回归”序列。如果函数 $f$ 是线性函数，我们称其为线性（自）回归。

\paragraph{向量情况。} 
为了简化符号，我们可以定义一个向量 $\vx \in \mathbb{R}^d$ 来收集序列中的 $d$ 个连续值：
\begin{equation}
    \vx_n \doteq [x_{n+d-1}, \ldots,  x_{n}]^\top, \quad \vx_n \in \mathbb{R}^d, \; n = 1, 2, 3, \ldots
\end{equation}
使用这种符号，递归关系 \eqref{eqn:recursive-d} 可以方便地写成：
\begin{equation}
    \vx_{n+1} = g(\vx_{n}) \; \in \mathbb{R}^d, \quad n =  1, 2, 3, \ldots
    \label{eqn:recursive-v}
\end{equation}
其中函数 $g(\cdot)$ 由 \eqref{eqn:recursive-d} 中的函数 $f$ 唯一定义，并且它以一个 $d$ 维向量作为输入。在不同的上下文中，这样的向量有时被称为“状态”或“标记”。注意，\eqref{eqn:recursive-d} 中的方程表示一个映射 $\mathbb{R}^d \rightarrow \mathbb{R}$，但这里的方程是 $g: \mathbb{R}^d \rightarrow \mathbb{R}^d$。

\paragraph{受控预测。}
我们也可以定义一个依赖于另一个可预测序列作为输入的可预测序列：
\begin{equation}
    \vx_{n+1} = f(\vx_{n}, \vu_n) \; \in \mathbb{R}^d, \quad n =  1, 2, 3, \ldots,
\label{eqn:recursive-control}
\end{equation}
其中 $\{\vu_n\}$ 且 $\vu_n \in \mathbb{R}^k$ 是一个（可计算的）可预测序列。换句话说，下一个向量 $\vx_{n+1} \in \mathbb{R}^d$ 既依赖于 $\vx_n \in \mathbb{R}^d$ 也依赖于 $\vu_n \in \mathbb{R}^k$。在控制理论的背景下，序列 $\{\vu_n\}$ 通常被称为“控制输入”，而 $\vx_n$ 被称为系统 \eqref{eqn:recursive-control} 的“状态”或“输出”。一个经典的例子是线性动力系统：
\begin{equation}
    \vx_{n+1} = \boldsymbol{A}\vx_n + \boldsymbol{B}\vu_n, \quad \boldsymbol{A} \in \mathbb{R}^{d\times d}, \boldsymbol{B} \in \mathbb{R}^{d\times k},
    \label{eqn:lineary-system} 
\end{equation}
这在控制理论中被广泛研究 \cite{Cal:Des}。

通常，控制输入是由状态 $\vx_n$ 本身的一个可计算函数给出的，例如：
\begin{equation}
    \vu_n = h(\vx_n), \quad n =  1, 2, 3, \ldots 
\end{equation}
因此，序列 $\{\vx_n\}$ 由两个可计算函数 $f$ 和 $h$ 的组合给出：
\begin{equation}
    \vx_{n+1} = f\big(\vx_{n}, h(\vx_n)\big), \quad n =  1, 2, 3, \ldots
    \label{eqn:recursive-closed-loop}
\end{equation}
这样，序列 $\{\vx_n\}$ 又变成了一个自回归可预测序列。当输入 $\vu_n$ 依赖于输出 $\vx_n$ 时，我们说产生的序列是由一个“闭环”系统 \eqref{eqn:recursive-closed-loop} 产生的。由于闭环系统不再依赖于任何外部输入，我们说这样的系统已经变得“自主”。它可以被视为自回归的一个特例。例如，在上述线性系统 \eqref{eqn:lineary-system} 中，如果我们选择 $\vu_n = \boldsymbol{F}\vx_n$，闭环系统变为：
\begin{equation}
        \vx_{n+1} = \boldsymbol{A}\vx_n + \boldsymbol{B}\vu_n = \boldsymbol{A}\vx_n + \boldsymbol{B}\boldsymbol{F}\vx_n = (\boldsymbol{A}+ \boldsymbol{B}\boldsymbol{F})\vx_n,
    \label{eqn:lineary-system-closed}
\end{equation}
这是一个线性自回归。

% \paragraph{近似预测。}
% 如果在上述方程 \eqref{eqn:recursive-control} 中输入序列 $\vu_n$ 不是可预测的，会怎样？一个极端的情况是它完全不可预测，例如：
% \begin{equation}
%     \vx_{n+1} = f(\vx_{n}, \boldsymbol{\epsilon}_n), \quad n =  1, 2, 3, \ldots
%     \label{eqn:recursive-noise}
% \end{equation}
% 其中 $\boldsymbol{\epsilon}_n \in \mathbb{R}^k$ 是一些随机（高斯）噪声。在这种情况下，我们说所定义的序列 $\{\vx_n\}$ 是“随机的”。显然，这样的序列不再完全可预测，至少不像上面那样在确定性意义上是可预测的。

% 当噪声的幅度 $\|\boldsymbol{\epsilon}_n\|$ 很小且函数 $f$ 不会显著放大噪声时，我们有
% \begin{equation}
%     \vx_{n+1} \approx f(\vx_{n}, \boldsymbol{0}) + \nabla_{\boldsymbol{\epsilon}}f(\vx_{n}, \boldsymbol{0})   \boldsymbol{\epsilon}_n, \quad n =  1, 2, 3, \ldots
%     \label{eqn:recursive-noise-small}
% \end{equation}
% 其中第二项的幅度 $\|\nabla_{\boldsymbol{\epsilon}}f(\vx_{n}, \boldsymbol{0})   \boldsymbol{\epsilon}_n\|$ 很小。因此，$ \vx_{n+1}$ 主要由第一项确定，其期望是可预测的：
% \begin{equation}
%     \mathbb{E}[\vx_{n+1}] \approx f(\vx_{n}, \boldsymbol{0}).
% \end{equation}
% 在这种情况下，我们说由 \eqref{eqn:recursive-noise-small} 定义的序列 $\{\vx_n\}$ 是“近似可预测的”。在实际应用中，我们经常处理这种近似可预测性，特别是在现实世界的数据中，测量噪声是常见的。需要注意的是，取期望本质上消除了随机噪声的影响，使得期望 $\{\mathbb{E}[\vx_n]\}$ 成为一个可预测的序列。这样的过程通常被称为“去噪”。\DP{这里“去噪”与第2/3章中的“去噪”之间的联系并不明显。例如，这里是对自回归序列而言，而后面是对单个观测而言。这与滤波也不是同一件事。我不太明白为什么需要这一段（至少在这里）。}

\paragraph{连续过程。}
可预测序列在连续情况下也有其自然的对应物。我们可以将它们称为可预测过程。类似于自然数序列，最简单的可预测连续过程是时间本身 $x=t$。

更一般地，我们说一个过程，记为 $\vx(t)$，是可预测的，如果在任何时间 $t$，该过程在 $t+\delta t$ 时的值，其中 $\delta t$ 是一个无穷小增量，由其在 $t$ 时的值确定。通常，值的变化 $\delta \vx(t)$ 是连续且平滑的。因此，$\delta \vx(t) = \vx(t + \delta t) - \vx(t)$ 是无穷小的。可预测过程通常由（多变量）微分方程描述：
\begin{equation}
    \dot{\vx}(t) = f(\vx(t)), \quad \vx \in \mathbb{R}^d. 
    \label{eqn:process}
\end{equation}

在系统理论中 \cite{Cal:Des,Sastry-Nonlinear}，上述方程也被称为状态空间模型。类似于离散情况，一个受控过程可以表示为：
\begin{equation}
    \dot{\vx}(t) = f(\vx(t), \vu(t) ), \quad \vx \in \mathbb{R}^d, \vu \in \mathbb{R}^k,
    \label{eqn:process-controlled}
\end{equation}
其中 $\vu(t)$ 是一个可计算的输入过程。

\begin{example}
    例如，在物理学中，牛顿第二定律描述了在力输入 $\boldsymbol{F}(t) \in \mathbb{R}^3$ 下，移动物体的轨迹 $\boldsymbol{x}(t) \in \mathbb{R}^3$ 的预测方法：
\begin{equation}
    m\ddot{\boldsymbol{x}}(t) = \boldsymbol{F}(t).
\end{equation}
当没有力 $\boldsymbol{F}(t) \equiv 0$ 时，上述定律简化为一个特殊情况，即牛顿第一定律：物体保持匀速直线运动：
\begin{equation}
   \ddot{\boldsymbol{x}}(t) = \boldsymbol{0} \; \Leftrightarrow \; \dot{\boldsymbol{x}}(t) = \boldsymbol{v}
\end{equation}
其中 $\boldsymbol{v} \in \mathbb{R}^3$ 是一个常数速度向量。
\end{example}

\subsection{低维度性}\label{sec:intro-low-dimensionality}
\paragraph{学习预测。}
现在假设你已经观察或被给予了许多序列片段：
\begin{equation}
    \{S_1, S_2, \ldots, S_i, \ldots, S_N\}
\end{equation}
所有这些片段都来自某个可预测序列 $\{x_n\}_{n=1}^\infty$。不失一般性，我们可以假设每个片段的长度为 $D \gg d$。因此，每个片段的形式为：
\begin{equation}
    S_i = [x_{j(i)}, x_{j(i)+1}, \ldots, x_{j(i)+D-1}]^\top \in \mathbb{R}^D
\end{equation}
其中 $j \in \mathbb{N}$。然后，你被给予一个新的片段 $S_t$，并被要求预测其未来的值。

这里的一个困难是，你通常不知道生成序列的函数 $f$ 和阶数 $d$：
\begin{equation}
    x_{n+d} = f(x_{n+d-1}, \ldots,  x_{n}).
\label{eqn:sequence-order-d}
\end{equation}
因此，希望以某种方式“学习” $f$ 和 $d$ 从给定的样本段 $S_1, S_2, \ldots, S_N$。因此，学习预测的核心任务是：
\begin{center}
{\em 给定许多可预测序列的采样段，如何有效地和高效地识别函数 $f$。}
\end{center}

\paragraph{可预测性和低维性。}
为了识别预测函数 $f$，我们可能会注意到任何可预测序列段的共同特征，例如由 \eqref{eqn:sequence-order-d} 给出的。如果我们取序列的一个长段，例如长度 $D \gg d$，并将其视为向量：
\begin{equation}
    \vx_i = [x_i, x_{i+1}, \ldots x_{i+D-1}]^\top \in \mathbb{R}^D.
\end{equation}
那么所有这样的向量集合 $\{\vx_i\}$ 远非随机，因此不可能占据整个 $\R^D$ 空间。相反，它们本质上最多有 $d$ 个自由度——给定任何 $\vx_i$ 的前 $d$ 个条目，其余条目的值是唯一确定的。换句话说，所有 $\{\vx_i\}$ 都位于一个 $d$ 维表面上。在数学中，这样的表面通常称为子流形，记为 $\mathcal{S} \subset \R^D$。

\begin{figure}[t]
\centering
\includegraphics[width=0.6\linewidth]{figures/two-dimensional plane in R10.pdf}
    \caption{十维空间中的二维子空间。}
    \label{fig:lowdimplane}
\end{figure}
在实践中，如果我们选择段长度 $D$ 足够大，那么从同一预测函数采样的所有段都位于一个具有内在维度 $d$ 的表面上，这远低于环境空间 $D$ 的维度。例如，如果序列由以下线性自回归给出：
\begin{equation}
    x_{n+2} = a\cdot x_{n+1} + b\cdot x_n,
    \label{eqn:sequence-2d}
\end{equation}
对于某些常数 $a, b \in \R$。如果我们从这样的序列中采样长度 $D =10$ 的段，那么所有样本都位于 $\R^{10}$ 中的一个二维平面或子空间上，如图 \ref{fig:lowdimplane} 所示。如果我们能够识别这个二维子空间，那么 \eqref{eqn:sequence-2d} 中的常数 $a$ 和 $b$ 可以完全确定。

很容易看出，如果预测函数 $f$ 是线性的，例如线性系统 \eqref{eqn:lineary-system} 和 \eqref{eqn:lineary-system-closed} 给出的情况，长段总是位于某个低维线性子空间上。识别预测函数在很大程度上等同于识别这个低维子空间，这个问题通常称为主成分分析。我们将在第 \ref{ch:classic} 章讨论这些经典模型和方法。

事实证明，对于一般的可预测序列也是如此：如果能够识别所有段样本都位于的低维表面，那么就可以识别相关的预测函数 $f$。\footnote{在温和条件下，低维表面和函数 $f$ 之间存在一一对应关系。这一事实已在系统识别等问题中得到利用，我们将在后面讨论。} 我们不能过分强调可预测序列段的这一性质的重要性：{\em 可预测序列的长段的所有样本都位于一个低维子流形上。} 正如我们将在本书中看到的，所有现代学习方法本质上都利用了这一性质，无论是隐式还是显式。 %\yima{添加一个图来说明表面或子流形。}

在现实世界中，我们观察到的数据通常不是来自单一的可预测序列。它们通常包含多个可预测序列的观测值。例如，一个视频序列可以包含多个移动物体。很容易看出，在这种情况下，数据位于多个低维线性子空间或非线性子流形的混合中，如图 \ref{fig:mixture-models} 所示。
\begin{figure}
    \centering
    \includegraphics[width=0.8\linewidth]{figures/mixture.pdf}
    \caption{分布在（正交）子空间（左）或子流形（右）的混合数据。}
    \label{fig:mixture-models}
\end{figure}

\paragraph{低维性的性质。}
当然，可预测序列中的时间相关性并不是数据低维性的唯一原因。例如，所有图像的空间是巨大的，但大部分空间由类似于无结构随机图像的图像组成，如图 \ref{fig:noise-image}(左) 所示。然而，自然图像和视频具有高度冗余性，因为所有像素值之间存在强烈的空域和时域相关性。这就是为什么即使部分被遮挡，我们也能轻松预测物体部分的外观，如图 \ref{fig:noise-image} 中间和右侧所示。因此，自然图像的分布具有非常低的内在维度（相对于图像的总像素数）。

\begin{figure}
    \centering
    
\includegraphics[height=0.30\linewidth]{figures/Gaussian-noise.png} \hspace{2mm} \includegraphics[height=0.30\linewidth]{figures/imageprocessing.png}
    \caption{随机噪声图像与自然图像及其损坏版本的对比。%\sdb{可能需要避免使用 Lena，因为 IEEE/PAMI-TC 等最近的动向}
    }
    \label{fig:noise-image}
\end{figure}

由于学习低维结构的任务的重要性和普遍性，书籍《高维数据分析与低维模型：原理、计算与应用》\cite{Wright-Ma-2022} 在开头就陈述道：“在高维空间中识别信号或数据的低维结构是最重要的问题之一，通过漫长的历史，它交织了许多工程和数学领域，如系统理论、信号处理、模式识别、机器学习和统计学。”

请注意，通过将观测数据点 $\vx$ 强制位于低维表面上，我们实际上使 $\vx$ 的条目之间非常依赖，并在某种意义上使条目非常“可预测”。例如，如果我们知道数据受限于 $\R^d$ 中的一个 $d$-维表面，那么除了预测之外，我们还可以对数据进行许多有用的操作：%\yima{可能需要一些这些性质的说明。}
\begin{itemize}
    \item \textbf{补全}：一般来说，给定一个典型样本 $\vx$ 的超过 $d$ 个条目，其剩余条目通常可以被 \emph{唯一} 确定。\footnote{预测成为这一性质的一个特例。}
    \item \textbf{去噪}：假设样本 $\vx$ 的条目受到 \emph{小} 噪声的干扰，可以通过将 $\vx$ 投影回表面来有效地去除这些噪声。
    \item \textbf{错误校正}：假设 $\vx$ 的一些 \emph{未知} 条目被任意损坏，可以有效地和高效地进行校正。
\end{itemize}
图 \ref{fig:low-dim-properties} 通过一个低维线性结构（二维平面上的一维直线）说明了这些性质。

\begin{figure}
    \centering
    \includegraphics[height=0.28\linewidth]{figures/Completion-low-dim.png}     \includegraphics[height=0.28\linewidth]{figures/Denoising-low-dim.png} \includegraphics[height=0.32\linewidth]{figures/Correction-low-dim.png} 
    \caption{低维（线性）结构的性质说明：它使补全（左）、去噪（中）和错误校正（右）成为可能。}
    \label{fig:low-dim-properties}
\end{figure}

实际上，在温和的条件下，上述性质可以推广到高维空间中的许多其他低维结构 \cite{Wright-Ma-2022}。有趣的是，正如我们将在本书中看到的，这些有用的性质，如补全和去噪，将启发有效的方法来学习这样的低维结构。

在上述内容中，为了简单起见，我们仅使用确定性情况来介绍可预测性和低维性的关键概念。因此，数据（或采样片段）精确地位于某些几何结构上，如子空间或曲面。在实践中，正如我们之前提到的，数据中总是存在一定程度的不确定性和随机性。在这种情况下，我们可以假设数据具有概率分布，其概率密度函数为 $p(\vx)$。我们说一个分布是“低维的”，如果其密度集中在某个低维的几何结构上，例如子空间、曲面或它们的混合，如图 \ref{fig:mixture-models} 所示。

我们上面的讨论已经引出了本书的主要论点：{\em 任何学习方法或智能系统都应该并且能够依赖于世界是可预测的，因此观察到的数据片段具有低维结构。} 剩下的问题是，如何通过有效且高效的可计算方法来正确学习这些低维结构。

\section{如何学习？}

\subsection{分析方法}
\label{sec:analytical}
%\paragraph{低维性的分析模型。}
即使一个预测函数易于计算，也不意味着从一些采样片段中学习这个函数是易于计算或可扩展的。当然，一种经典的保证问题易于计算或易于高效求解的方法是明确假设我们所处理的低维结构族。历史上，由于计算和数据的限制，简单而理想化的分析模型总是首先被研究，因为它们通常提供高效的闭式解或数值解。此外，它们可以为更一般的问题提供见解，并且通常已经为重要但有限的情况提供了有用的解决方案。在计算资源稀缺的旧时代，允许高效闭式解或数值解的分析模型是唯一可以实现的案例。{\em 线性结构}成为首先被彻底研究的模型类别。

例如，可以说最简单的情况是假设数据在高维空间中的单个低维子空间附近分布。或者，等价地，可以假设数据根据一个几乎退化的低维高斯分布。从有限数量的（噪声）样本中识别这样的子空间或高斯分布，就是经典的主成分分析（PCA）问题，为此类模型开发了有效的算法 \cite{JolliffeI2002}。可以逐步增加模型族的复杂性和表达能力。例如，可以假设数据分布在某些低维成分（子空间或低维高斯分布）的混合上，如独立成分分析（ICA） \cite{Ans-1985}、字典学习（DL）、广义主成分分析（GPCA） \cite{Vidal-GPCA} 或近年来在压缩感知等领域广泛研究的更一般的稀疏低维模型 \cite{Wright-Ma-2022}。

所有这些分析模型家族的核心研究问题始终是如何在每个家族中找到最紧凑的模型，该模型最适合给定的数据。下面，我们简要介绍这些经典分析模型，但将更系统的研究留到第 \ref{ch:classic} 章。理论上，这些分析模型为我们提供了关于低维结构的几何和统计性质的巨大见解。它们通常为我们提供了闭式解或高效且可扩展的算法，这些算法对于可以很好地用此类模型近似的数据非常有用。更重要的是，对于更一般的问题，它们为我们提供了一种一般的感觉，即识别低维结构的问题有多容易或困难，以及解决此类问题的基本思想是什么。

%Hence, in this book, we choose to introduce and study these somewhat idealistic structures first in Chapter \ref{ch:classic}, not only because we can derive efficient solutions for them with rigorous theoretical guarantees, but also the basic ideas of their solutions carry over to the more general (nonlinear) structures and distributions that we will study in later chapters.

\subsubsection{线性动态系统}
\label{sec:linear-systems}

\paragraph{维纳滤波器。}

正如我们在第 \ref{sec:predictability} 节中讨论过的，智能的主要任务之一是学习在观察序列中什么是可预测的。可能最简单的可预测序列或信号类是通过线性时不变（LTI）过程生成的：
\begin{equation}
    x[n] = h[n]*z[n] + \epsilon[n], 
    \label{eqn:Wiener-model}
\end{equation}
其中 $z$ 是输入，$h$ 是脉冲响应函数。\footnote{通常假设 $h$ 具有某些良好的结构，例如有限长度或带状频谱。} 这里 $\epsilon[n]$ 是观察中的加性噪声。问题是，给定输入过程 $\{z[n]\}$ 和输出过程 $\{x[n]\}$ 的观察，如何找到最优的 $h[n]$ 使得 $\hat x[n] = h[n]*z[n]$ 以最优的方式预测 $x[n]$。通常，我们通过最小均方误差（MMSE）来衡量预测的好坏：
\begin{equation}
    \min_{h} \mathbb{E} \big[\epsilon[n]^2\big] = \mathbb{E} \big[\|x[n] - h[n]*z[n]\|_2^2\big].
\end{equation}
最优解 $h[n]$ 被称为（去噪）滤波器。诺伯特·维纳（Norbert Wiener），同为控制论运动的发起人，在 1940 年代研究了这个问题，并给出了一个优雅的闭式解，称为维纳滤波器 \cite{Wiener-1942,Wiener-1949}。这成为信号处理领域最重要的结果之一。

\paragraph{卡尔曼滤波器。} 
动态过程的去噪或滤波思想后来由鲁道夫·卡尔曼（Rudolph Kalman）在 1960 年代扩展到由状态空间模型（有限维）描述的线性时不变系统：
\begin{equation}
    \z[n] = \boldsymbol{A} \z[n-1] + \boldsymbol{B}\boldsymbol{u}[n] + \boldsymbol{\epsilon}[n]. 
    \label{eqn:linear-state-space}
\end{equation}
问题是，如何从形式为 $\x[n] = \boldsymbol{C} \z[n] + \boldsymbol{w}[n]$ 的噪声观察中估计系统的状态 $\z[n]$，
\label{eqn:Kalman-model}
其中 $\boldsymbol{w}$ 是某种（白）噪声。最小化 MMSE 类预测误差的最优因果\footnote{这意味着估计只能使用当前时间步 $n$ 之前的观察。卡尔曼滤波器总是因果的，而维纳滤波器不一定。} 状态估计器由所谓的卡尔曼滤波器 \cite{kalman1960new} 给出。这是现代控制理论的基石之一，因为它允许我们从动态系统的噪声观察中估计其状态。然后，可以随后引入一个（线性）状态反馈，例如形式为 $\boldsymbol{u}[n] = \boldsymbol{F} \hat{\boldsymbol{z}}[n]$，并使闭环系统完全自主，正如我们在方程 \eqref{eqn:recursive-closed-loop} 中看到的那样。

\paragraph{线性动力系统的识别。}
为了推导卡尔曼滤波器，假设系统参数 $(\boldsymbol{A}, \boldsymbol{B}, \boldsymbol{C})$ 是已知的。如果这些参数没有事先给出，那么这将是一个更具挑战性的问题，称为{\em 系统识别}：如何从输入序列 $\{\boldsymbol{u}[n]\}$ 和观测序列 $\{\x[n]\}$ 的多个样本中{\em 学习}参数 $(\boldsymbol{A}, \boldsymbol{B}, \boldsymbol{C})$。这是系统理论中的一个经典问题。如果系统是线性的，可以证明输入和输出序列 $\{\boldsymbol{u}[n], \x[n]\}$ 会共同位于某个低维子空间上\footnote{该子空间的维度与状态空间模型 \eqref{eqn:linear-state-space} 的阶数相同。}。因此，识别问题本质上等同于识别这个低维子空间 \cite{OverscheeP1996,Liu-2009-CDC,Liu-2010-SIAM}。

请注意，上述问题有两个共同点：首先，假设（无噪声）序列或信号是由一组显式的参数模型生成的；其次，这些模型本质上都是线性的。因此，概念上，设 $\x_o$ 是一个随机变量，其“真实”分布支持在一个低维线性子空间上，记为 $S$。在很大程度上，维纳滤波器和卡尔曼滤波器都试图从其噪声观测中估计这样的 $\x_o$：
\begin{equation}
    \x = \x_o + \boldsymbol{\epsilon}, \quad \x_o \sim S, 
\end{equation}
其中 $\boldsymbol{\epsilon}$ 通常是随机高斯噪声（或过程）。因此，本质上，它们的解决方案都依赖于识别一个最佳拟合观测（噪声）数据的低维线性子空间。然后，通过将数据投影到这个子空间上，可以获得最优的去噪操作，所有这些操作都是封闭形式的。

%系统理论、信号/图像处理、统计学习和机器学习中的平行发展：如何在高维空间中有效地和高效地学习一个更一般的低维分布？线性系统的去噪（维纳滤波器和卡尔曼滤波器）的历史。经验

\subsubsection{线性和混合线性模型}
\label{sec:PCA-ICA}
\paragraph{主成分分析。}
从经典信号处理和系统识别中的上述问题，我们看到所有这些问题背后的主要任务是从噪声观测中学习一个{\em 单个}低维线性子空间。数学上，我们可以将这种结构建模为：
\begin{equation}
    \x = \boldsymbol{u}_1 z_1 + \boldsymbol{u}_2 z_2 + \cdots + \boldsymbol{u}_d z_d + \boldsymbol{\epsilon} =  \boldsymbol{U} \z + \boldsymbol{\epsilon}, \quad \boldsymbol{U} \in \mathbb{R}^{D\times d}
    \label{eqn:PCA-model}
\end{equation}
其中 $\boldsymbol{\epsilon} \in \mathbb{R}^D$ 是一些小的随机噪声。图 \ref{fig:PCA} 说明了具有两个主成分的这种分布。
\begin{figure}
    \centering
    \includegraphics[width=0.5\linewidth]{figures/PCA.png}
    \caption{具有两个主成分的分布。}
    \label{fig:PCA}
\end{figure}
问题是根据 $\x$ 的多个样本找到子空间基 $\boldsymbol{U}$。估计子空间 $\boldsymbol{U}$ 的典型方法是通过最小化噪声的方差，也称为最小均方误差（MMSE）：
\begin{equation}
    \min_{\boldsymbol{U}} \mathbb{E}\big[\|\boldsymbol{\epsilon}\|_2^2\big] = \mathbb{E}\big[\|\x - \boldsymbol{U} \z\|_2^2\big].
\end{equation}
请注意，这本质上是一个去噪任务：一旦正确找到基 $\boldsymbol{U}$，我们可以通过将其投影到由 $\boldsymbol{U}$ 张成的低维子空间上来去噪噪声样本 $\x$：
\begin{equation}
\x \rightarrow \hat \x = \boldsymbol{U}\boldsymbol{U}^\top \x. 
\end{equation}
如果噪声很小，并且我们正确地学习了低维子空间 $\boldsymbol{U}$，我们应该期望 $\x \approx \hat \x$。也就是说，PCA 是自编码的一个特例：
\begin{equation}
    \x   \xrightarrow{\hspace{2mm} \boldsymbol{U}^\top\hspace{2mm}} \z  \xrightarrow{\hspace{2mm} \boldsymbol{U} \hspace{2mm}} \hat \x.
       \label{eqn:auto-encoding-PCA}
\end{equation}
这里，由于数据结构简单，编码器 $\mathcal{E}$ 和解码器 $\mathcal{D}$ 都成为简单的线性算子（投影和提升）。

这是统计学中的一个经典问题，称为主成分分析（Principal Component Analysis, PCA）。它最早由Pearson在1901年研究 \cite{Pearson1901}，后来由Hotelling在1933年独立研究 \cite{Hotelling1933}。这个主题在经典著作 \cite{Jolliffe1986,JolliffeI2002} 中得到了系统的总结。

此外，人们可能会明确假设数据 $\x$ 服从单个低维高斯分布：
\begin{equation}
    \x \sim \mathcal{N}(\boldsymbol{0}, \boldsymbol{U}\boldsymbol{U}^\top + \sigma \I), \quad \boldsymbol{U} \in \mathbb{R}^{D\times d},
\end{equation}
这等价于假设上述PCA模型 \eqref{eqn:PCA-model} 中的 $\z$ 是标准正态分布。这被称为概率PCA \cite{TippingM1999}。

在本书中，我们将在第 \ref{ch:classic} 章重新审视PCA，从学习低维分布的角度出发。我们的目标是使用这个简单而理想化的模型来传达一些学习低维分布紧凑表示的基本思想，包括通过去噪和自编码来压缩表示的重要概念。

\paragraph{独立成分分析。}
独立成分分析（ICA）最初由 \cite{Ans-1985} 作为学习良好表示的经典模型提出。事实上，它最初被提出作为我们记忆的简单数学模型。ICA模型与上述PCA模型 \eqref{eqn:PCA-model} 有着看似相似的形式，假设观测随机变量 $\x$ 是多个独立成分 $z_i$ 的线性叠加：
\begin{equation}
    \x = \boldsymbol{u}_1 z_1 + \boldsymbol{u}_2 z_2 + \cdots + \boldsymbol{u}_d z_d  + \boldsymbol{\epsilon} =  \boldsymbol{U} \z + \boldsymbol{\epsilon}.
    \label{eqn:ICA-model}
\end{equation}
然而，这里的成分 $z_i$ 被假设为独立的非高斯变量。例如，一个流行的选择是
\begin{equation}
    z_i = \sigma_i \cdot w_i, \quad \sigma_i \sim B(1,p),
    \label{eqn:ICA-modes}
\end{equation}
其中 $\sigma_i$ 是伯努利随机变量，$w_i$ 可以是常数值或另一个随机变量，例如高斯分布。\footnote{即使 $w$ 是高斯分布，$\sigma w$ 也不再是高斯变量！} ICA问题试图从观测样本 $\x$ 中识别出 $\boldsymbol{U}$ 和 $\z$。图 \ref{fig:ICA-PCA} 说明了ICA和PCA之间的区别。

\begin{figure}
    \centering
    \includegraphics[width=0.7\linewidth]{figures/PCA_ICA.png}
    \caption{PCA（左）与ICA（右）。}
    \label{fig:ICA-PCA}
\end{figure}

尽管从 $\z$ 到 $\x$ 的（解码）映射在学习 $\boldsymbol{U}$ 和 $\z$ 后看起来是线性和简单的，但从 $\x$ 到 $\z$ 的（编码）映射可能很复杂，可能无法用简单的线性映射表示。因此，ICA通常给出如下形式的自编码：
\begin{equation}
    \x   \xrightarrow{\hspace{2mm} \mathcal{E}\hspace{2mm}} \z  \xrightarrow{\hspace{2mm} \boldsymbol{U} \hspace{2mm}} \hat \x.
       \label{eqn:auto-encoding-ICA}
\end{equation}
因此，与PCA不同，ICA的分析和求解稍微困难一些。在1990年代，像Erkki Oja和Aapo Hyv\"{a}rinen \cite{hyvarinen-1997,Hyvrinen-2000} 等人对ICA做出了重要的理论和算法贡献。在第 \ref{ch:classic} 章中，我们将研究并给出ICA的解决方案，其中编码映射 $\mathcal{E}$ 将变得清晰。

\paragraph{稀疏结构和压缩感知。}
如前所述，如果 \eqref{eqn:ICA-modes} 中的 $p$ 非常小，那么任何成分非零的概率都很小。在这种情况下，我们说 $\x$ 是稀疏生成的，并且它集中在维度为 $k = p \cdot d$ 的线性子空间上。因此，在某种程度上，我们可以将上述ICA模型扩展到更广泛的低维结构，即稀疏模型。

一个$k$-稀疏模型被定义为所有{\em $k$-稀疏向量}的集合：
\begin{equation}
    \Z = \{\z \in \mathbb{R}^n \mid \|\z\|_0 \le k\},
\end{equation}
其中$\| \cdot \|_0$是$\ell^0$-范数，它计算向量$\z$中非零元素的数量。也就是说，$\Z$是所有与坐标轴对齐的$k$-维子空间的并集，如图\ref{fig:mixture-models}左所示。经典信号处理和统计学中的一个重要问题是，如何从线性观测中恢复一个稀疏向量：
\begin{equation}
    \x = \boldsymbol{A} \z, \quad \boldsymbol{A} \in \mathbb{R}^{m\times n}
    \label{eqn:sparse-model}
\end{equation}
其中$\boldsymbol{A}$是已知的，但通常$m < n$。这个看似简单的 问题实际上是非常困难的，计算上是NP-hard的，甚至难以近似（详见参考文献\cite{Wright-Ma-2022}）。

尽管有非常丰富和悠久的研究历史，可以追溯到18世纪\cite{Boscovichca1750}，但直到20世纪初，还没有一种能够证明其有效性的算法来解决这类问题，尽管在20世纪60年代到90年代之间提出了许多启发式算法。其中一些在实践中非常有效，但没有严格的理论依据。直到2000年代初，一些著名的数学家，包括David Donoho、Emmanuel Cand\`{e}s和Terence Tao\cite{donoho2005neighborly,Candes2005,CandesE2005-IT}建立了严格的理论框架，使我们能够精确地描述在哪些条件下稀疏恢复问题可以有效地解决，例如通过凸$\ell^1$最小化：
\begin{equation}
    \min \|\z\|_1 \quad \mbox{subject to} \quad \x = \boldsymbol{A}\z,
\end{equation}
其中$\|\cdot \|_1$是向量的稀疏促进$\ell^1$范数。这个问题的任何解本质上都给出了一个稀疏编码映射：
\begin{equation}
    \x   \xrightarrow{\hspace{2mm} \mathcal{E} \hspace{2mm}}  \z.
       \label{eqn:decoding-sparse}
\end{equation}
我们将在第\ref{ch:classic}章中简要介绍这种算法，即映射，这将揭示稀疏编码与深度学习之间有趣的本质联系。\footnote{尽管早在2010年就注意到了稀疏编码算法与深度网络算法之间的相似性\cite{gregor2010learning}。}

事实证明，$\ell^1$最小化成功的条件非常广泛。成功恢复所需的最小测量次数$m$仅与数据的内在维度$k$成比例。这现在被称为{\em 压缩感知}现象\cite{CandesE2006-ICM}。此外，这种现象不仅限于稀疏结构，还适用于非常广泛的低维结构，如低秩矩阵等。这些结果从根本上改变了我们对高维空间中低维结构恢复的理解。这种高维数据分析的戏剧性逆转甚至被David Donoho\cite{DonohoD2000}称为{\em ``维度的祝福''}，与高维问题中典型的悲观信念“维度的诅咒”形成鲜明对比。这一连贯且完整的成果体系已被系统地整理在参考文献\cite{Wright-Ma-2022}中。

从计算的角度来看，这一新框架的重要性无法被高估。它从根本上改变了我们对一类以前被认为难以处理的问题的看法。它使我们能够开发出极其高效的算法，这些算法随着问题维度的增加而优雅地扩展，从而使得稀疏恢复问题从：
\begin{equation}
    \mbox{\textbf{不可处理}} \;
   \Longrightarrow \; \mbox{\textbf{可处理}} \; \Longrightarrow \; 
   \mbox{\textbf{可扩展}}.
\end{equation}
这些算法还提供了关于其正确性的严格理论保证，这些保证基于对数据和计算的精确要求。这种方法的严格性和精确性几乎与实践中的深度神经网络相反，后者在很大程度上是经验性的。然而，尽管它们的风格和标准看似截然不同，但我们现在知道这两种方法都共享一个共同的目标：{\em 在高维空间中追求低维结构}。

\paragraph{字典学习。}
概念上，比稀疏编码问题 \eqref{eqn:sparse-model} 更难的问题是当观测矩阵 $\boldsymbol{A}$ 甚至在事先都不知道，而需要从一组观测中学习时：
\begin{equation}
    \X = \boldsymbol{A} \Z, \quad \boldsymbol{A} \in \mathbb{R}^{m\times n}.
    \label{eqn:dictionary-learning}
\end{equation}
这被称为{\em 字典学习}问题，可以视为对之前讨论的ICA问题 \eqref{eqn:ICA-model} 的推广。换句话说，给定数据 $\X$ 的分布为标准稀疏分布 $\Z$ 在线性变换 $\boldsymbol{A}$ 下的图像，我们希望学习 $\boldsymbol{A}$ 和其“逆”映射 $\mathcal{E}$，使得我们获得一个自编码器：
\begin{equation}
    \X   \xrightarrow{\hspace{2mm} \mathcal{E} \hspace{2mm}}  \Z \xrightarrow{\hspace{2mm} \boldsymbol{A} \hspace{2mm}} \X?
       \label{eqn:decoding-DL}
\end{equation}

可以看出，PCA、ICA 和字典学习的共同点在于它们都假设数据的分布围绕低维线性或混合线性结构。它们都需要从可能有噪声的分布样本中学习线性结构的（全局或局部）基。在第 \ref{ch:classic} 章中，我们将研究如何通过这些经典模型来识别低维结构。特别是，我们将看到一个有趣且重要的事实：所有这些低维（分段）线性模型都可以通过同一种快速算法，即{\em 幂迭代} \cite{Zhai-2020}，有效地学习。尽管上述线性或混合线性模型对于大多数现实世界的数据来说过于简单或理想化，但理解这些模型是理解更一般的低维分布的重要第一步。

\subsubsection{一般分布}\label{sec:denoising-intro}

现实世界的数据，如图像、视频和音频，其分布过于复杂，无法用上述理想化的线性模型或高斯过程来建模。我们通常不知道它们是从什么参数模型族生成的。{\footnote{尽管历史上曾尝试开发这些数据的分析模型，例如图像数据的随机场或随机过程 \cite{Mumford-1999}，如我们在前一节中讨论的。}}在实践中，我们通常只有许多来自其分布的样本——经验分布。显然，在这种情况下，我们通常不能期望有任何闭式解来解决它们的低维结构，更不用说由此产生的去噪算子了。{\footnote{与PCA、维纳滤波器和卡尔曼滤波器的情况不同。}}因此，我们需要为这些经验分布开发一个更一般的解决方案，不一定以闭式形式存在，但至少可以高效计算。如果我们正确地做到了这一点，上述线性模型的解决方案应该成为其特例。

在20世纪50年代，统计学家开始对从任意分布中抽取的数据去噪问题感兴趣。设 $\x_o$ 是一个具有概率密度函数 $p_o(\cdot)$ 的随机变量。因此，如果我们观察到 $\x_o$ 的一个噪声版本：
\begin{equation}
    \x = \x_o + \boldsymbol{\epsilon}, 
\end{equation}
其中 $\boldsymbol{\epsilon}$ 是一个高斯噪声 $\boldsymbol{\epsilon} \sim \mathcal{N}(\boldsymbol{0}, \sigma^2 \boldsymbol{I})$。设 $p(\cdot)$ 是 $\x$ 的概率密度函数，\footnote{即，$p(\x) = \int_{-\infty}^{\infty} \phi_{\sigma}(\x - \z)p_o(\z)~d\z$，其中 $\phi_{\sigma}$ 是高斯分布 $\mathcal{N}(\boldsymbol{0}, \sigma^2 \boldsymbol{I})$ 的密度函数。}令人惊讶的是，给定 $\x$ 的条件下 $\x_o$ 的后验期望可以通过一个优雅的公式计算，这个公式被称为Tweedie公式 \cite{Robbins1956AnEB}：\footnote{Herbert Robbins 将这个公式的功劳归功于 Maurice Kenneth Tweedie，这是他们个人通信中的内容。}
\begin{equation}
    \hat \x_o = \mathbb{E}[\x_o\mid \x] = \x + \sigma^2 \nabla \log p(\x).
\end{equation}
从这个公式可以看出，函数 $\nabla \log p(\x)$ 在这里扮演了一个非常特殊的角色，用于去噪观测值 $\x$。噪声 $\boldsymbol{\epsilon}$ 可以显式地估计为
\begin{equation}
    \hat{\boldsymbol{\epsilon}} = \x - \hat \x_o = \sigma^2 \nabla \log p(\x),
\end{equation}
我们只需要知道 $\x$ 的分布 $p(\cdot)$，而不需要知道 $\x_o$ 的真实分布 $p_o(\cdot)$。这个结果的一个重要推论是，如果我们向任何分布添加高斯噪声，去噪过程可以通过某种方式获取函数 $\nabla \log p(\x)$ 来轻松完成。

\begin{figure}
    \centering
    %\includegraphics[width=0.45\linewidth]{figures/Score.png}
    \includegraphics[width=0.8\linewidth]{figures/Density-gradient.png}
    \caption{密度 $p(\x)$ 的分布的得分函数 $\nabla \log p(\x)$ 的几何解释。得分函数生成的操作将分布推向密度更高的区域。目标是，通过某种紧凑性的度量，使结果分布更加“压缩”。例如，熵 $H(\x) = - \int p(\x) \log p(\x)~d \x$ 或分布 $R(\x)$ 的编码率（通过某种有损编码方案）会变得更小。}
    \label{fig:score-function}
\end{figure}

实际上，这个函数有一个非常直观的信息论和几何解释。注意，在信息论中 $-\log p(\x)$ 通常对应于编码 $\x$ 所需的比特数\footnote{至少在离散变量的情况下，我们将在第 \ref{ch:compression} 章中更详细地解释这一点。}。梯度 $\nabla \log p(\x)$ 指向密度更高的方向。因此，如果 $\x$ 向该方向移动，编码 $\x$ 所需的比特数会减少。因此，算子 $\nabla \log p(\x)$ 的总体效果是将分布推向密度更高的区域，从而减少结果分布的总体编码长度/率，使其更加“压缩”。实际上，可以正式证明，通过这种操作得到的分布的熵会减少。图 \ref{fig:score-function} 说明了与得分函数相关的这种操作的几何意义。

由于这是一个如此重要和有用的结果，它在许多不同的背景和领域中被重新发现和使用。例如，几年后它被 \cite{Miyasawa61} 重新发现。在21世纪初，函数 $\nabla \log p(\x)$ 在学习一般分布的背景下再次被重新发现，并被 Aapo Hyv\"{a}rinen \cite{hyvarinen05a} 命名为“得分函数”。但其与（经验贝叶斯）去噪的联系很快被 \cite{Vincent2011} 认识到。Eero Simoncelli 的团队 \cite{Raphan10} 对其他测量分布（超出高斯噪声）进行了推广，后来应用于图像去噪 \cite{Kadkhodaie21a,ho2020denoising}。

我们将稍后在本章和第 \ref{ch:compression} 章中讨论，这样一个看似简单的去噪和压缩概念如何导致了一种非常统一且强大的方法，用于在高维空间中学习一般的低维分布，包括自然图像的分布。

\subsection{经验方法}
在实践中，对于许多重要的现实世界数据，如图像、声音和文本，很难用理想化的线性或混合线性模型来建模。例如，在图像处理和计算机视觉领域，长期以来一直试图通过解析方法来建模自然图像的分布。菲尔兹奖得主大卫·芒福德（David Mumford）在 1990 年代花费了大量精力尝试理解和建模自然图像的统计特性 \cite{Mumford1996TheSD}。他和他的学生，包括宋春朱，从统计物理学中汲取灵感和技术，提出了许多用于自然图像分布的统计和随机模型 \cite{Zhu-Entropy-1997,Zhu1997LearningGP,Zhu1997Prior,Huang-Mumford,Mumford-1999,Lee-Mumford}。然而，这些解析模型在生成与自然图像非常相似的样本方面遇到了有限的成功。显然，对于像图像这样的现实世界数据，我们需要开发更强大和统一的方法来追求它们更一般的低维结构。

因此，历史上提出了许多经验模型来建模重要的现实世界数据，包括图像和文本。这些模型通常从生物神经系统的特点中汲取灵感，因为动物或人类的大脑似乎能够极其高效和有效地处理这些数据。

\subsubsection{经典人工神经网络}
\paragraph{人工神经元。}

\begin{figure}[t]
    \centering
\includegraphics[width=0.55\linewidth]{figures/neuron.png} \hspace{3mm}   
\includegraphics[width=0.40\linewidth]{figures/Artificial_neuron.png}
    \caption{第一个模仿神经元（左）处理信号的人工神经元（右）的数学模型。}
    \label{fig:neuron}
\end{figure}

受大脑神经系统的启发，第一个人工神经元的数学模型\footnote{称为线性阈值单元或感知器。} 由沃伦·麦卡洛克\footnote{当时芝加哥大学的精神病学教授} 和沃尔特·皮茨于 1943 年提出 \cite{McCulloch-Pitts}。它描述了输入 $x_i$ 和输出 $o_j$ 之间的关系：
\begin{equation}
    o_j = \varphi\Big( \sum_i w_{ji}x_i\Big),  
\end{equation}
其中 $\varphi(\cdot)$ 是某种非线性激活函数，通常由阈值函数建模。该模型如图 \ref{fig:neuron} 所示。我们可以看到，这种形式已经具有现代深度神经网络基本单元的主要特征。该模型是从观察我们神经系统中单个神经元的工作方式得出的。然而，人们并不确切知道一组这样的神经元想要实现和执行什么功能。在更技术的层面上，他们也不确定应该使用什么非线性激活函数 $\varphi(\cdot)$。因此，历史上提出了许多变体。\footnote{阶跃函数、硬或软阈值、修正线性单元（ReLU）、sigmoid 等。}

\begin{figure}[t]
\centering
\includegraphics[width=0.85\linewidth]{figures/single-deep.png}
    \caption{A network with one single hidden layer (left) versus a deep network (right).  }
    \label{fig:single-deep}
\end{figure}
\paragraph{人工神经网络。}
在20世纪50年代，Frank Rosenblatt是第一个构建具有这种人工神经元网络的机器的人，如图\ref{fig:perceptron}所示。这台机器被称为Mark I Perceptron，它由一个输入层、一个输出层和一个由512个人工神经元组成的单个隐藏层组成，如图\ref{fig:perceptron}左所示，这与图\ref{fig:single-deep}左所示的类似。它被设计用于分类光学图像中的字母。然而，单层网络的容量有限，只能学习线性可分的模式。在Marvin Minsky和Seymour Papert于1969年出版的《感知器：计算几何导论》一书中，证明了Mark I Perceptron的单层架构无法学习异或（XOR）函数。这一结果极大地抑制了人们对人工神经网络的兴趣，尽管后来证明多层网络能够学习异或函数\cite{Rumelhart1986}。事实上，如图\ref{fig:single-deep}右所示，由这些简单神经元组成的多层网络（足够大）可以模拟任何有限状态机，甚至通用图灵机。\footnote{不要混淆神经网络在原则上能够做什么与是否能够有效地学习实现某些期望功能的神经网络。} 然而，随后，在20世纪70年代，人工神经网络的研究进入了第一个寒冬。

\begin{figure}
    \centering
    \includegraphics[width=0.45\linewidth]{figures/visu-large.jpg}
    \hspace{2mm} \includegraphics[width=0.45\linewidth]{figures/Original-Mark-I-perceptron-as-seen-in-its-operators-manual-20.ppm.png}
    \caption{Frank Rosenblatt在20世纪50年代末开发的Mark I Perceptron机器。}
    \label{fig:perceptron}
\end{figure}

\paragraph{卷积神经网络。}
20世纪50年代和60年代，像Mark I Perceptron这样的早期人工神经网络实验结果令人失望，这表明仅仅以多层感知器（MLP）的方式连接神经元可能还不够。为了构建有效且高效的网络，我们需要理解网络中的神经元需要共同实现什么目的或功能，以便它们应该以某种特殊方式组织和学习。再次，机器智能的研究转向从动物的神经系统中寻找灵感。

众所周知，我们的大脑大部分用于处理视觉信息。在20世纪50年代和60年代，David Hubel和Torsten Wiesel系统地研究了猫的视觉皮层。他们发现视觉皮层包含不同类型细胞（称为简单细胞和复杂细胞），这些细胞对不同方向和位置的视觉刺激敏感\cite{Hubel-Wiesel-1959}。Hubel和Wiesel因这一开创性发现而获得1981年诺贝尔生理学或医学奖。

在人工神经网络方面，Hubel和Wiesel的工作启发了Kunihiko Fukushima，他在1980年设计了“新认知器”网络，该网络由模仿视觉皮层中生物神经元的人工神经元组成\cite{Fukushima1980NeocognitronAS}。这被称为第一个{\em 卷积神经网络}（CNN），其架构如图\ref{fig:neocognitron}所示。与感知器不同，新认知器具有多个隐藏层，可以被视为深度网络，如图\ref{fig:single-deep}右所示。

同样受到猫的视觉皮层中神经元工作方式的启发，他也是第一个引入{\em 矩形线性单元}（ReLU）的人：
\begin{equation}
    \varphi(x) = \max\{0, x\} = \left\{ 
    \begin{matrix}
        x, \quad \mbox{如果} \, x > 0,\\
        0, \quad \mbox{如果} \, x \le 0,
    \end{matrix}
    \right.
\end{equation}
作为激活函数 $\varphi(\cdot)$ 在1969年 \cite{Fukushima-1969}。但直到近年来，ReLU 才成为现代深度（卷积）神经网络中广泛使用的激活函数。一旦我们解释了深度网络试图实现的主要操作：压缩，我们就会从这本书中了解到为什么这是一个好的选择。
%\yima{这里放一张卷积神经网络或LeNet的图片...}

\begin{figure}
    \centering
    \includegraphics[width=0.6\linewidth]{figures/neocognitron.png}
    \caption{卷积神经网络的起源：Kunihiko Fukushim 在1980年提出的 Neocognitron。请注意，卷积层和池化层的交错尝试模仿在猫的视觉皮层中发现的简单细胞和复杂细胞的功能。}
    \label{fig:neocognitron}
\end{figure}

CNN 类型的网络在1980年代继续发展，许多不同的变体被引入和研究。然而，尽管深度网络具有显著的能力，以及受神经科学启发的改进架构，训练这样的深度网络以完成诸如图像分类等实际任务仍然极其困难。如何让网络工作取决于许多无法解释的启发式方法和技巧，这确实限制了神经网络的吸引力和适用性。1989年左右，Yann LeCun 成功使用{\em 反向传播}（BP）来学习一个用于识别手写数字的深度卷积神经网络 \cite{LeCun-1989}，这成为了一个重大突破，后来被称为 LeNet（见图 \ref{fig:LeNet-5}）。经过几年的持续发展，他的坚持得到了回报：LeNet 的性能最终在1990年代末变得足够好，可以用于实际应用 \cite{LeCun-1998}：它被美国邮政局用于识别手写数字（用于邮政编码）。LeNet 被认为是所有现代深度神经网络的“原型”网络，例如 AlexNet 和 ResNet，我们将在后面讨论。由于这项工作，Yann LeCun 获得了2018年图灵奖。\footnote{与另外两位深度网络的先驱 Yoshua Bengio 和 Geoffrey Hinton 一起。}

\begin{figure}
    \centering
\includegraphics[width=0.95\linewidth]{figures/LeNet-5.png}
    \caption{Yann LeCun 在1989年设计的 LeNet-5 网络。}
    \label{fig:LeNet-5}
\end{figure}

\paragraph{反向传播。}
在历史上，深度神经网络的命运似乎与它们如何能够被轻松高效地训练紧密相关。反向传播（BP）就是为了这个目的而引入的。我们知道，多层感知器可以表示为一系列线性映射和非线性激活函数的组合，如下所示：
\begin{equation}
h(\bm W_1, \ldots, \bm W_L) = f^L(\bm W_Lf^{L-1}(\bm W_{L-1} \cdots f^2(\bm W_2f^1(\bm W_1\x)))).
\end{equation}
为了训练网络权重 $\{\bm W_l\}_{l=1}^L$，以减少基于梯度下降算法的预测/分类误差，我们需要评估梯度 ${\partial h}/{\partial \bm W_l}$。从微积分中的{\em 链式法则}可知，这种类型的函数的梯度可以高效地计算，后来被称为反向传播（BP）。详见附录 \ref{app:optimization} 的详细描述。反向传播技术在 1960 年代和 1970 年代就已经被人们在最优控制和动态规划等领域所熟知和实践。例如，它出现在 Paul Werbos 博士的 1974 年博士论文中 \cite{Werbos-1974, Werbos1994TheRO}。1986 年，David Rumelhart 等人首次将反向传播应用于训练多层感知器（MLP）网络 \cite{Rumelhart1986}。从此以后，BP 逐渐流行起来，因为它提供了一种{\em 可扩展}的算法来学习大型深度神经网络。\footnote{因为它可以在促进并行和分布式计算的计算平台上高效实现。}现在，它已成为当今训练深度神经网络的几乎主导技术。然而，人们普遍认为，自然界并不通过反向传播来学习，因为这种机制对于自然界来说仍然过于昂贵，无法实现物理化。\footnote{正如我们之前讨论的，自然界几乎普遍通过闭环反馈来学习纠正错误。}这显然为未来留下了巨大的改进空间，我们将在后面进一步讨论。

然而，尽管在 1980 年代取得了上述算法进展和有希望的实践，但在 1980 年代和 1990 年代，训练深度神经网络仍然对计算系统来说非常挑剔和昂贵。在 1990 年代末，支持向量机（SVM） \cite{SVM-1995} 已经非常流行，因为它们被视为神经网络在分类等任务上的更好替代方案。\footnote{事实上，解决分类问题的类似思想可以追溯到 Thomas Cover 的博士论文工作，该工作在 1964 年被浓缩并发表在一篇论文中 \cite{Cover-1964}。}有两个主要原因：首先，SVM 基于一种称为 Vapnik--Chervonenkis（VC）理论的严格统计学习框架；其次，它基于凸优化 \cite{BoydVa04} 导致了相当高效的算法。SVM 的兴起在 2000 年初左右给神经网络的研究带来了第二个寒冬。

\paragraph{压缩自动编码器。}
在20世纪80年代末和90年代，人工神经网络已经被用于学习高维数据（如图像）的低维表示。已经证明，神经网络可以用于从数据中学习主成分分析（PCA），而不是使用第\ref{sec:PCA-ICA}节中讨论的经典方法。在20世纪80年代末，人们还争论说，由于其能够建模非线性变换的能力，神经网络被建议用于学习具有非线性分布的数据的低维表示。与线性PCA的情况类似，人们可以尝试同时学习一个非线性降维编码器$f$和一个解码器$g$，每个都由一个深度神经网络建模\cite{Rumelhart1986,Kramer1991NonlinearPC}：
\begin{equation}
    \X   \xrightarrow{\hspace{2mm} f \hspace{2mm}} \Z  \xrightarrow{\hspace{2mm} g \hspace{2mm}} \hat \X.
       \label{eqn:auto-encoding-deep-networks}
\end{equation}
通过强制解码数据$\hat \X$与原始数据$\X$一致，例如通过最小化一种均方误差（MMSE）类型的重构误差\footnote{尽管已知MMSE类型的误差对于具有复杂非线性结构的图像数据存在问题。正如我们稍后将讨论的，最近在生成方法中的大部分工作，包括GANs，都是为了找到原始数据$\X$和再生数据$\hat \X$之间更好的距离函数的替代品。}：
\begin{equation}
    \min_{f,g} \big\|\X - \hat \X\big\|_2^2 = \big\|\X - g(f( \X))\big\|_2^2,
\end{equation}
可以从数据$\X$本身学习一个自动编码器。

但是，我们如何保证这样的自动编码器确实捕获了$\X$中的真实低维结构，而不是给出一个冗余的表示呢？例如，我们可以简单地选择$f$和$g$为恒等映射，并且$\Z = \X$。因此，为了确保自动编码器是有价值的，人们希望得到的表示在某种可计算的复杂性度量下是压缩的。1993年，Geoffrey Hinton及其同事提出了使用编码长度作为这样的度量，因此自动编码器的目标变为找到一个最小化编码长度的表示\cite{Hinton-1993}。这项工作还建立了最小描述长度原理\cite{Rissanen-1978}与自由（亥姆霍兹）能量最小化之间的基本联系。Hinton小组的后续工作\cite{Hinton504}实证表明，这样的自动编码器能够学习现实世界图像的有意义的低维表示。在深度网络变得流行之前，Pierre Baldi在2011年进行了一项更全面的自动编码器综述\cite{Baldi2011}。我们将在第\ref{sec:unifying-approach}节中更详细地讨论复杂性度量和自动编码器，并在第\ref{ch:consistent}章中以更统一的观点系统研究压缩自动编码器。

\begin{figure}
    \centering
    \includegraphics[width=0.8\linewidth]{figures/Comparison_image_neural_networks.svg.png}
    \caption{LeNet \cite{LeCun-1989}与AlexNet \cite{Krizhevsky2012imagenet}的架构对比。}
    \label{fig:LeNet-AlexNet}
\end{figure}

\subsubsection{现代深度神经网络}
在近30年的时间里，从20世纪80年代到2010年代，对于机器学习和机器智能的研究，神经网络并没有被主流认真考虑。早期（深度）神经网络，如LeNet，在小规模分类问题（如识别数字）上表现出色。然而，网络的设计和实践非常经验化，当时可用的数据集较小，反向传播算法对于当时的计算机来说是一个巨大的计算负担。这些因素导致了对神经网络的兴趣缺乏，进展停滞不前，只有少数研究人员继续研究它。

\paragraph{分类与识别。}
事实证明，深度神经网络的巨大潜力只有在有足够的数据和计算能力时才能释放。快进到2010年代，更大的数据集如ImageNet变得可用，GPU变得足够强大，使得BP（反向传播）对于比LeNet更大的网络也变得更加可行。大约在2012年，一个名为AlexNet的深度卷积神经网络因其在ImageNet数据集上的表现超过了现有的分类方法而引起了关注 \cite{Krizhevsky2012imagenet}。\footnote{实际上，在此之前，深度网络已经在语音识别任务上展示了最先进的性能。但直到它们在图像分类上的成功才受到如此多的关注。} 图 \ref{fig:LeNet-AlexNet} 显示了AlexNet和LeNet之间的比较。AlexNet与LeNet有许多共同的特征，只是它更大，并且采用了ReLU作为非线性激活函数，而不是LeNet中使用的Sigmoid函数。部分由于这项工作的影响力，Geoffrey Hinton获得了2018年的图灵奖。

这一早期的成功激发了机器智能社区在接下来的几年中探索新的网络设计变体和改进。特别是，人们发现，网络越大越深，在图像分类等任务中的性能就越好。许多深度网络架构被尝试、测试并普及。一些值得注意的包括VGG \cite{Simonyan15}、GoogLeNet \cite{Szegedy2014GoingDW}、ResNet \cite{He2016-lc}以及最近的Transformers \cite{vaswani2017attention}等。尽管在经验性能方面取得了快速进展，但这些经验发现的架构缺乏理论解释，包括它们之间的关系（如果有）。本书的一个目的是揭示所有这些网络可能共同服务的共同目标，以及为什么它们共享某些共同特征，包括多层线性算子与非线性激活函数的交错（见第 \ref{ch:representation} 章）。

\paragraph{强化学习。}
深度网络的早期成功主要是在监督学习设置中的分类任务，如语音识别和图像识别。后来，深度网络被Demis Hassabis领导的DeepMind团队采用，用于学习游戏中的决策或控制策略。在这种情况下，深度网络用于建模最优决策/控制策略或相关的最优值函数，如图 \ref{fig:Alpha-Go} 所示。这些网络参数是基于当前策略下游戏的成功或失败所返回的奖励进行逐步优化的。\footnote{例如，基于反向传播（BP）。} 这种学习方法通常被称为{\em 强化学习} \cite{Sutton-Barto}，起源于20世纪60年代末控制系统的实践 \cite{Waltz1965AHA,Mendel1970ReinforcementlearningCA}。它的早期根源可以追溯到Richard Bellman的{\em 动态规划} \cite{Bellman-DP} 和Marvin Minsky的{\em 试错学习} \cite{Minsky-1954} 在20世纪50年代的工作。

从实现的角度来看，深度网络和强化学习的结合证明是非常强大的：深度网络可以用于近似难以用解析方法建模的真实世界环境中的控制策略和值函数。这种实践最终导致了DeepMind公司开发的AlphaGo系统，在2016年击败了顶级人类选手李世石，在围棋比赛中，然后在2017年击败了世界冠军柯洁。\footnote{1996年，IBM的Deep Blue系统通过击败俄罗斯国际象棋大师加里·卡斯帕罗夫而创造了历史。它主要使用传统的机器学习技术，如树搜索和剪枝，这些技术在可扩展性和对更复杂的游戏（如围棋）的成功方面并不那么有效，超过20年。}

AlphaGo的成功对计算社会来说是一个巨大的惊喜，因为普遍认为搜索的状态空间太大，以至于无法找到任何有效的解决方案，无论是从计算还是样本量的角度来看。它成功的唯一合理解释是，围棋游戏的最优值/策略函数中必须存在非常好的结构。它们的内在维度并不高，可以通过神经网络很好地近似，可以从不是特别多的样本中学习。

\begin{figure}
    \centering
    \includegraphics[width=0.5\linewidth]{figures/Policy-Value.png}
    \caption{AlphaGo：使用深度神经网络来建模围棋游戏的最优策略或最优值函数。}
    \label{fig:Alpha-Go}
\end{figure}

\paragraph{生成与预测。}
人们可以将2010年代早期的深度网络实践视为更多地关注从数据$\X$中提取相关信息并将其编码为特定任务的表示$\Z$（例如，在分类任务中，$\Z$表示类别标签）：
\begin{equation}
    \X   \xrightarrow{\hspace{2mm} f \hspace{2mm}} \Z.
       \label{eqn:encoding-deep-networks}
\end{equation}
在这种设置下，需要学习的映射$f$不需要保留关于$\X$的大部分分布信息，只需要保留特定任务所需的充分统计量。例如，$\X$中的样本$\x$可能是一张苹果的图像，由$f$映射到其类别标签$\z =$ ``苹果''。{\em 信息瓶颈框架} \cite{Tishby-ITW2015}于2015年提出，用于分析深度网络在这种上下文中的作用。

然而，在许多现代情况下，例如那些所谓的大型基础模型，人们经常需要解码$\Z$以在一定程度上恢复相应的$\X$：
\begin{equation}
    \Z   \xrightarrow{\hspace{2mm} g  \hspace{2mm}} \hat \X.
       \label{eqn:decoding-deep-networks}
\end{equation}
由于$\X$通常表示从外部世界观察到的数据，一个好的解码器将允许我们模拟或预测世界中发生的事情。例如，在“文本到图像”或“文本到视频”任务中，$\z$通常表示描述所需图像$\x$内容的文本。解码器应该能够生成一个$\hat \x$，其内容与$\x$相同。例如，给定一个对象类别$\z = $ ``苹果''，解码器$g$应该生成一个看起来像苹果的图像$\hat \x$，尽管不一定与原始$\x$完全相同。

\paragraph{通过判别方法生成。}
为了使生成的图像$\hat \X$与真实的自然图像$\X$相似，我们需要能够评估并最小化某个距离：
\begin{equation}
    \min d(\X, \hat \X).
\end{equation}
事实证明，大多数理论驱动的距离在高维空间中但具有低内在维度的分布上计算和优化都非常困难，如果不是不可能的话。\footnote{即使给定了$\X$的参数分布族，距离也经常变得病态或未定义，对于低维支撑的分布。更糟糕的是，所选的族可能无法很好地近似感兴趣的真正分布。}

2007年，宋春朱的学生朱若文可能对早期尝试通过分析方法建模和生成自然图像感到失望（之前讨论过），决定尝试一种截然不同的方法。在2007年发表于CVPR的一篇论文中\cite{Tu-2007}，他首次提出可以通过判别式方法学习图像的生成模型。这个想法很简单：如果难以评估距离$d(\X, \hat \X)$，可以尝试学习一个判别器$d$来区分$\hat \X$和$\X$：
\begin{equation}
    \Z   \xrightarrow{\hspace{2mm} g  \hspace{2mm}} \hat \X, \X \xrightarrow{\hspace{2mm} d  \hspace{2mm}} \boldsymbol{0}, \boldsymbol{1},
       \label{eqn:gan-networks}
\end{equation}
其中$\boldsymbol{0}, \boldsymbol{1}$表示图像是否生成。直觉上，我们越能区分$\hat \X$和$\X$，它们可能越接近。

朱的工作\cite{Tu-2007}首次展示了通过判别式方法学习生成模型的可行性。然而，该工作采用了传统的图像生成和分布分类方法（如提升法），这些方法速度慢且难以实现。2012年后，深度神经网络在图像分类中变得非常流行。2014年，Ian Goodfellow及其同事再次提出使用判别式方法生成自然图像\cite{Goodfellow-2014}。他们建议使用深度神经网络来建模生成器$g$和判别器$d$。此外，他们提出了通过极小极大博弈学习生成器$g$和判别器$d$：
\begin{equation}
    \min_g \max_d \ell(\X, \hat \X),
\end{equation}
其中$\ell(\cdot)$是与分类相关的某种自然损失函数。换句话说，判别器$d$试图最大化其在区分$\X$和$\hat \X$方面的成功，而生成器$g$则试图相反。因此，这种方法被称为{\em 生成对抗网络}（GANs）。研究表明，一旦在大量数据集上训练，GANs确实可以生成逼真的图像。部分由于这项工作的影响力，Yoshua Bengio获得了2018年的图灵奖。

判别式方法似乎是一种巧妙地绕过分布学习基本困难的方法。然而，严格来说，这种方法并没有完全解决这一基本困难。\cite{Goodfellow-2014}表明，通过选择合适的损失函数，极小极大公式在数学上等价于最小化{\em Jensen-Shannon距离}（参见\cite{Cover-Thomas}）。这在高维空间中对于两个低维分布来说是一个困难的问题。因此，在实践中，GANs通常依赖于许多启发式方法和工程技巧，并且经常遭受诸如{\em 模式崩溃}等不稳定问题的影响。\footnote{尽管如此，这种极小极大公式提供了一种实用的距离近似方法。它简化了实现并避免了直接计算距离时的一些陷阱。}总体而言，关于如何改进GANs的实践，缺乏理论指导。

\paragraph{通过去噪和扩散生成。}
2015年，GAN被引入并变得流行后不久，Surya Ganguli及其学生意识到并建议，可以通过一个由深度网络建模的迭代去噪过程来学习一般分布，例如自然图像的分布\cite{Sohl-Dickstein2015}。他们的方法受到了William Feller在1949年研究的特殊高斯和二项过程的启发\cite{Feller1949OnTT}。

\begin{figure}[t]
    \centering
    \includegraphics[width=\linewidth]{figures/diffusion_pipeline-2.pdf}
    \caption{从高斯分布开始，通过迭代去噪过程收敛到任意数据分布的示意图。}
    \label{fig:diffusion}
\end{figure}

不久之后，基于得分函数的去噪算子（在第 \ref{sec:denoising-intro} 节中简要介绍）被证明更加通用，并统一了去噪和扩散过程及算法 \cite{song2019,song2020score,ho2020denoising}。图 \ref{fig:diffusion} 展示了一个将通用高斯分布转换为任意经验分布的过程，通过执行一系列迭代去噪（或压缩）操作。到目前为止，去噪（和扩散）已经取代了生成对抗网络（GANs），成为学习图像和视频分布的主流方法，从而产生了流行的商业图像生成引擎，如 Midjourney 和 Stability.ai。

在第 \ref{ch:compression} 章中，我们将系统地介绍和研究用于学习一般低维分布的去噪和扩散方法。

\section{一种统一的方法}\label{sec:unifying-approach}
到目前为止，我们已经简要介绍了机器智能的主要目标和历史，以及与其相关的许多重要思想和方法。近年来，在深度神经网络的经验成功之后，人们付出了巨大的努力来开发理论框架，以帮助我们理解所有经验设计的深度神经网络，无论是某些看似必要的组件（例如，dropout、归一化、注意力等）还是它们的整体行为（例如，双下降、神经坍缩等）。

部分受到这一动机的启发，本书旨在实现几个重要且具有挑战性的目标：
\begin{itemize}
    \item 开发一个理论框架，使我们能够对深度神经网络进行严格的数学解释。
    \item 确保学习到的数据分布的正确性，并与学习到的表示保持一致。
    \item 证明该框架可以导致高性能的架构，并指导进一步的实践改进。
\end{itemize}
在过去的几年中，有越来越多的证据表明，通过利用早期讨论的经典分析低维模型的理论和解决方案（在第 \ref{ch:classic} 章中更详细地讨论），以及整合来自几个相关领域的基本思想，即信息/编码理论、控制/博弈理论和优化，这些目标确实可以实现。本书旨在系统地介绍这一新方法。

\subsection{学习简约表示}
\label{sec:computational-approach-compression}
任何学习任务可能的一个必要条件是，感兴趣的序列必须是可计算的，至少在艾伦·图灵的意义上是如此 \cite{Turing-1936}。也就是说，一个序列可以通过在典型计算机上运行的程序来计算。\footnote{确实存在不可计算的定义良好的序列。} 除了可计算性之外，我们还要求计算是可处理的。\footnote{我们不需要考虑那些计算复杂度不可处理的事物，例如随着序列长度或维度的增长而呈指数增长的事物。} 也就是说，学习和计算序列的计算成本（空间和时间）不应随长度呈指数增长。此外，正如我们在自然界（以及现代机器智能实践中）所看到的，对于大多数实际任务，智能系统需要从高维空间中的大量数据中学习可预测的内容，例如从视觉、声音和触觉中学习。因此，对于智能，我们不需要考虑所有可计算和可处理的序列或结构。我们应该只关注那些允许其学习和计算算法具有可扩展实现的可预测序列和结构：
\begin{equation}
\mbox{\textbf{可计算}} \;
   \Longrightarrow \; \mbox{\textbf{可处理}} \; \Longrightarrow \; 
   \mbox{\textbf{可扩展}}。
\end{equation}

这是因为智能生物用来学习有用信息的任何算法都必须是{\em 可扩展的}。更具体地说，算法的计算复杂度最好能够优雅地扩展，通常线性或甚至亚线性地扩展到数据的大小和维度。在技术层面上，这要求算法所依赖的操作只能利用可以从数据中高效计算的预言信息。更具体地说，当维度较高且规模较大时，唯一能够计算的预言信息要么是关于数据的一阶几何信息\footnote{例如，用线性子空间局部近似非线性结构，或计算目标函数的梯度。}，要么是关于数据或其特征的二阶统计信息\footnote{例如，数据或其特征的协方差或相关性。}。

本书的主要目标是在一个理论和计算框架内，系统地开发具有这种可扩展预言和操作的高效和有效的解决方案或算法，以从采样数据中{\em 学习}低维结构，并随后学习预测函数。

\paragraph{通过压缩追求低维性。}
从我们在第 \ref{sec:predictability} 节中给出的序列示例中，很明显有些序列易于建模和计算，而其他序列则更难。显然，序列的计算成本取决于预测函数 $f$ 的复杂程度。回归次数 $d$ 越高，计算成本就越高。$f$ 可以是一个简单的线性函数，也可以是一个可以任意复杂指定和计算的非线性函数。

合理地认为，如果一个序列更难，无论我们选择什么度量标准来指定和计算它，那么从其采样片段中学习它也会更加困难。然而，对于任何给定的可预测序列，实际上有多种，通常是无限多种，方法来指定它。例如，对于一个简单的序列 $x_{n+1} = a x_{n}$，我们也可以用 $x_{n+1} = a x_n + b x_{n-1} - b x_{n-1}$ 来定义相同的序列。
因此，如果我们能够为任何给定的可计算序列开发一个客观和严格的“复杂度”概念，那将非常有用。

俄罗斯数学家安德烈·柯尔莫哥洛夫是最早为任何可计算序列定义复杂度的人之一。\footnote{许多人都对这一序列复杂度的概念做出了贡献，其中最著名的是雷·索洛蒙诺夫和格雷格·柴廷。据信，这三个人独立地发展和研究了算法信息论，雷·索洛蒙诺夫在 1960 年，安德烈·柯尔莫哥洛夫在 1965 年 \cite{Kolmogorov1998OnTO}，格雷戈里·柴廷在 1966 年左右 \cite{Chaitin-1966}。}他建议，在所有能够计算相同序列的程序中，我们可以使用最短程序的长度作为其复杂度的度量。这个想法与著名的“奥卡姆剃刀”原则非常一致：{\em 在所有能够解释相同观察结果的理论中，总是选择最简单的那一个。}更精确地说，设 $p$ 表示一个可以在通用计算机 $\mathcal{U}$ 上生成序列 $S$ 的计算程序。序列 $S$ 的柯尔莫哥洛夫复杂度定义为：
\begin{equation}
    K(S) = \min_{p\,:\, \mathcal{U}(p) = S} L(p). 
\end{equation}
因此，序列的复杂度通过我们能够以多“简约”的方式指定或计算它来衡量。这种序列复杂度或简约性的定义具有重要的概念意义，并且具有许多有趣的性质。历史上，它激发了许多关于计算理论的深刻研究，特别是在算法信息论领域。

最短程序的长度可以被视为所考虑序列的终极压缩，提供了我们通过学习序列的正确生成机制所获得多少收益的定量度量。然而，尽管它在理论上具有重要意义，但 Kolmogorov 复杂度通常不是可计算的函数 \cite{Cover-Thomas}（甚至难以准确近似）。因此，这种复杂度的度量在实际应用中几乎没有用处。它既不能告诉我们学习给定序列的难度，也不能告诉我们我们学习得有多好。

%简约或紧凑性的度量：许多（大多等价）方法来衡量数据的紧凑性：Kolmogorov 复杂度（不可计算）、数据分布（支撑集）的维度和体积、最小描述长度、率失真等。

\paragraph{可计算的简约度量。}
因此，为了实际目的，我们需要一个对从同一预测函数生成的序列进行高效计算的复杂度度量。\footnote{请注意，在实践中，我们通常关心学习预测函数 $f$，而不是 $f$ 生成的任何特定序列。} 需要注意的是，Kolmogorov 复杂度不可计算的部分原因是其定义是非构造性的。

因此，为了引入一个可计算的复杂度度量，我们可以采取更构造性的方法，正如 Claude Shannon 通过信息论框架所倡导的那样 \cite{Shannon-1948,Cover-Thomas}。\footnote{这在过去 80 多年中成功指导了通信行业的工程实践。} 本质上，通过假设序列 $S$ 从概率分布 $p(S)$ 中抽取，所谓的 {\em 熵}：
\begin{equation}
    H(S) \doteq -\int p(S) \log p(S)~dS
    \label{eqn:entropy-definition}
\end{equation}
给出了其复杂度的自然度量。这个度量也有一个自然的解释，即编码此类序列所需的平均二进制位数，我们将在第 \ref{ch:compression} 章中看到。

为了说明这一观点的主要思想，我们考虑一个由预测函数$f$生成的大量长序列片段：
\begin{equation}
    \{S_1, S_2, \ldots, S_i, \ldots, S_N\} \subset \mathbb{R}^D,
\end{equation}
注意，不失一般性，这里我们假设所有序列的长度都是$D$。因此，每个序列可以被视为$\mathbb{R}^D$中的一个向量。其次，我们可能引入一种编码方案（带有代码簿），记为$\mathcal E$，将每个片段$S_i$转换为一个唯一的二进制位流$\mathcal{E}(S_i)$。\footnote{最简单的编码方案是用$\epsilon$-球填充所有片段所张成的空间。然后，我们对所有球进行编号。每个序列被编码为其最近球的（二进制）编号。} 因此，每个片段可以从其对应的位流中恢复\footnote{如果编码方案是有损的，则可以达到一定的精度。}。然后，预测函数$f$的复杂度可以被评估为所有序列的平均编码长度，即编码率：
\begin{equation}
   R(f \mid \mathcal E) = \mathbb{E}[L(\mathcal{E}(S))] \approx \frac{1}{N}\sum_{i=1}^N L(\mathcal{E}(S_i)). 
   \label{eqn:coding-rate}
\end{equation}
显然，如果使用不同的编码方案（和代码簿），编码率度量将发生变化。在实践中，我们对片段分布的低维结构了解得越多，设计的代码簿就越高效。给定两个片段的编码方案$\mathcal{E}_1$和$\mathcal{E}_2$，如果编码率的差异为正：
\begin{equation}
   R(f \mid \mathcal E_1) -  R(f \mid \mathcal E_2) > 0, 
\end{equation}
我们可以说编码方案$\mathcal{E}_2$更好。这种差异可以被视为衡量$\mathcal{E}_2$相对于$\mathcal{E}_1$关于数据分布的更多信息量的度量。在很大程度上，学习$f$的目标等同于找到最高效的编码方案，以最小化编码率：
\begin{equation}
   \min_{\mathcal{E}} R(f \mid \mathcal E). 
\end{equation}
正如我们在第\ref{ch:compression}章中将看到的，可实现的最小率与熵$H(S)$\eqref{eqn:entropy-definition}密切相关。

\begin{figure}
    \centering
    \includegraphics[width=0.5\linewidth]{figures/Coding-schemes.png}
    \caption{两种编码方案的比较。假设数据的真实分布围绕绿色球体。可以使用由所有蓝色球体组成的代码簿对事件进行编码；也可以使用仅由绿色球体组成的代码簿对事件进行编码。显然，第二种编码方案将导致更低的编码长度/率。}
    \label{fig:coding-schemes}
\end{figure}

\begin{remark}\label{rem:computable-complexity}
    {使用显式编码方案衡量数据复杂性的视角激发了几个学习目标，这些目标旨在改进可计算性，包括1968年提出的最小消息长度（MML）\cite{WallaceC1999}和1978年提出的最小描述长度（MDL）\cite{Rissanen-1978,HansenM2001}。这些目标通常会计算编码方案$\mathcal{E}$本身（包括其代码簿）的编码长度，以及感兴趣的数据$S$：$L(\mathcal E(S)) + L(\mathcal E)$。然而，如果目标是学习一个有限大小的代码簿并将其应用于大量序列，那么代码簿的摊销成本可以忽略不计，因为当$N$变得很大时，$$\frac{1}{N}\Big( L(\mathcal{E}) + \sum_{i=1}^N L(\mathcal{E}(S_i))\Big) \approx \frac{1}{N}\sum_{i=1}^N L(\mathcal{E}(S_i))$$。}
\end{remark}

再次，可以将得到的最优编码方案视为实现对观测数据的最佳压缩。一般来说，与 Kolmogorov 复杂度相比，任何编码方案给出的编码长度总是更大：

\begin{equation}
    K(S) < L( \mathcal E(S)).
\end{equation}

因此，最小化编码速率/长度本质上是将不可计算的 Kolmogorov 复杂度的上界最小化。

\subsection{学习有信息量的表示}
%\paragraph{通过信息增益获得结构化表示。}
需要注意的是，如果目标仅仅是简单地压缩给定数据，只是为了压缩，那么在理论上，接近 Kolmogorov 复杂度的最优编码将变得几乎随机或无结构 \cite{Chaitin-1966}。\footnote{因为任何具有结构的编码都可以进一步压缩。} 然而，我们学习预测函数 $f$ 的真正目的是为了在未来预测中轻松重复使用它。因此，虽然压缩允许我们识别数据中的低维分布，但我们希望以结构化和有序的方式编码该分布，以便得到的表示非常有信息量且易于使用。正如我们在第 \ref{ch:representation} 章中将展示的那样，通过选择自然的编码方案族和基于该编码方案的自然的信息增益度量，可以精确地引入最终表示中所需的结构。这是采取明确和构造性编码方法的另一个优势。

\begin{figure}
    \centering
    \includegraphics[width=0.9\linewidth]{figures/Expansion.png}
    \caption{将数据分布转换为更具有信息量的结构表示。}
    \label{fig:expansion}
\end{figure}

正如我们将在本书中看到的那样，这种方法为学习低维结构的良好表示提供了一个强大的计算框架。在许多重要的实际情况下，编码长度函数可以高效地计算或准确地近似。在某些良性情况下，我们甚至可以获得闭式公式，例如子空间和高斯分布（见第 \ref{ch:compression} 章）。

此外，这种可计算和可量化的方法导致了一种原则性的方法，自然揭示了深度网络在学习过程中的作用。正如我们在第 \ref{ch:representation} 章中将系统推导的那样，它们的层正是试图以增量的方式实现优化感兴趣的目标函数的操作。从这个角度来看，深度网络的作用可以精确地解释为模拟某种迭代算法，例如梯度下降，以优化信息增益的目标。由此产生的深度架构的层可以赋予精确的统计和几何解释，即执行增量压缩编码和解码操作。因此，由此推导出的深度网络成为透明的“白盒”。

% \textbf{注释：} 将可预测性转化为低维性：此类数据的一个普遍的数学性质是它们的分布总是具有非常低的内在维度，因此在通常的高维空间中高度可压缩。

% % 促进低维性的额外相关性，例如许多由偏微分方程描述的物理定律，例如麦克斯韦方程组描述的电磁场和纳维-斯托克斯方程描述的流体动力学。

% 紧致性的度量：有许多（大多等价）方法来度量数据的紧致性：Kolmogorov 复杂度（不可计算）、数据分布（支撑集）的维度和体积、最小描述长度、率失真等。

% 因此，学习可预测性继承的低维结构的一种有效方法是通过某些“压缩”操作（我们将在后面精确化），这些操作旨在最小化这些度量。

\subsection{学习一致的表示}

%\subsection{确保一致性}
\label{sec:consistency}
为了总结我们到目前为止的讨论，让我们将数据表示为：
\begin{equation}
    \boldsymbol{X} = \{S_1, S_2, \ldots, S_i, \ldots, S_N\} \subset \mathbb{R}^D,
\end{equation}
并设 $\boldsymbol{Z} = \mathcal{E}(\X)$ 是通过某个编码器 $\mathcal{E}$ 对 $\X$ 的编码：
\begin{equation}
    \X  \xrightarrow{\hspace{2mm} \mathcal{E}\hspace{2mm}} \Z.
    \label{eqn:open-loop}
\end{equation}
在机器学习的背景下，$\Z$ 通常被称为“特征”或“潜在表示”。请注意，如果不了解 $\X$ 的底层分布，我们不知道哪个编码器 $\mathcal{E}$ 能够保留关于 $\X$ 分布的最有用信息。在实践中，人们通常从尝试某种特定的紧凑编码方案开始，这种方案对特定任务非常有效。特别是，他们会尝试学习一个编码器，该编码器优化了学习表示的某种（经验）简约度量：
\begin{equation}
    \min \rho(\Z). 
\end{equation}

\begin{example}
例如，图像分类就是这种情况：我们将同一类中的所有图像分配给一个单一的代码，将不同类中的图像分配给不同的代码，例如：
\begin{equation}
  \x \mapsto \z \in \{  [1, 0, 0, \ldots , 0, 0], \;  [0, 1, 0 \ldots, 0, 0], \; \ldots, \;  [0,0,0, \ldots, 0, 1].\}
  \label{eqn:class-labels}
\end{equation}
现在，分类器 $f(\cdot)$ 可以建模为一个函数，该函数预测给定 $\x$ 属于每个 $k$ 类的概率：$\hat{\z} = f(\x) \in \mathbb{R}^k$。然后，分类器的“好坏”可以通过所谓的{\em 交叉熵}来衡量：\footnote{交叉熵可以被视为 $\z$ 的真实分布与预测 $\hat\z$ 的分布之间的距离度量。它也可以被视为使用 $\hat \z$ 的最优码本对 $\z$ 进行编码时的预期编码长度。当 $\z$ 和 $\hat \z$ 具有相同的分布时，交叉熵达到最小值。}
\begin{equation}
    L(\hat{\z}, \z) = \sum_{i=1}^k - \z(i) \log \hat{\z}(i),
\end{equation}
其中 $\z(i)$ 表示向量 $\z$ 的第 $k$ 个元素。正如深度网络的早期实践所表明的 \cite{Krizhevsky2012imagenet}，如果给定足够的数据，这种编码方案通常可以通过深度网络表示，并通过优化交叉熵以端到端的方式学习。
\end{example}

交叉熵损失 $L(\hat \z, \z)$ 可以被视为与特定编码方案族相关的简约度量 $\rho(\z)$ 的一种特殊度量，这些编码方案适用于分类。然而，这种编码显然是{\em 非常有损的}。学习到的 $\z$ 除了类类型外，不包含关于 $\x$ 的任何其他信息。例如，通过将图像分配给（表示类标签“苹果”的代码），我们不再仅从标签中知道原始图像中的具体苹果类型。

当然，另一个极端是要求编码方案是{\em 无损的}。也就是说，$\x$ 和其代码 $\z$ 之间存在一对一的映射。然而，正如我们在第 \ref{ch:compression} 章中将看到的，除非 $\x$ 的值是离散的，否则无损编码（或压缩）是不切实际的。对于连续随机变量，我们只能考虑有损编码方案，以便数据的编码长度是有限的。也就是说，我们只对数据进行编码，直到达到某个预定的精度。正如我们在第 \ref{ch:compression} 章中将详细阐述的，有损编码不仅是实际选择，它在从分布的有限样本中学习潜在的连续分布方面起着根本作用。

对于许多学习目的，我们希望特征 $\z$，尽管是{\em 有损的}，但比仅保留其类别类型的信息更多。在本书中，我们将介绍一种基于更广泛编码方案族的简约性更一般的度量——混合子空间编码。这个族具有在一定精度内近似任意现实世界分布的能力。正如我们在第 \ref{ch:representation} 章中将看到的，这样的度量不仅能够保留关于 $\X$ 分布的大部分信息，还能促进学习表示 $\Z$ 的某些理想的几何和统计结构。

\paragraph{双向自编码以保持一致性。}
在一个更广泛的学习背景下，压缩编码方案 $\mathcal{E}$ 的主要目标是识别数据 $\X$ 中的低维结构，以便它们可以用于预测原始数据空间中的事物。这需要学习的编码方案 $\mathcal{E}$ 允许一个高效的解码方案，记为 $\mathcal D$。它将 $\Z$（通常称为潜在表示）映射回数据空间：
\begin{equation}
    \X   \xrightarrow{\hspace{2mm} \mathcal{E}\hspace{2mm}} \Z  \xrightarrow{\hspace{2mm} \mathcal{D} \hspace{2mm}} \hat \X.
       \label{eqn:auto-encoding}
\end{equation}
一般来说，我们称这样的编码和解码对 $(\mathcal{E}, \mathcal{D})$ 为{\em 自编码}。图 \ref{fig:autoencoder} 说明了这样一个自编码器的过程。
\begin{figure}
    \centering
    \includegraphics[width=0.7\linewidth]{figures/Autoencoder.jpg}
    \caption{自编码器架构的说明。}
    \label{fig:autoencoder}
\end{figure}

通常，我们希望解码尽可能接近编码的“逆”，使得从 $\Z$ 解码的数据 $\hat \X$ 在某种程度上与原始数据 $\X$ 相似。\footnote{我们将在后面更精确地说明相似的含义。} 如果是这样，我们就能从 $\Z$ 恢复或预测原始数据空间中发生的事情。在这种情况下，我们说对 $(\mathcal{E}, \mathcal{D})$ 给出了一致的自编码。对于大多数实际用途，不仅需要这样的解码存在，而且还需要能够高效实现和物理实现。例如，如果 $\x$ 是实值变量，为了在有限状态机上实现任何解码方案，需要量化，我们将在第 \ref{ch:compression} 章中更详细地解释这一点。因此，一般来说，应该期望可实现的编码和解码方案是必然有损的。一个中心问题是如何学习一个紧凑（有损）的表示 $\Z$，使得它可以用于很好地预测 $\X$。

一般来说，正如我们将看到的，编码器和解码器都可以通过深度网络建模和实现，并通过求解以下形式的优化问题来学习：
\begin{equation}
   \min \, d( \X, \hat \X) + \rho(\Z), 
   \label{eqn:auto-encoding-objective}
\end{equation}
其中 $d(\cdot, \cdot)$ 是某种距离函数，促进 $\X$ 和 $\hat \X$ 之间的相似性\footnote{根据距离函数 $d$ 的选择，可以是样本级或分布级的相似性。}，而 $\rho(\Z)$ 是某种度量，促进 $\Z$ 的简约性和信息丰富性。经典的主成分分析（PCA）\cite{JolliffeI2002} 是一致自编码的一个典型例子，我们将在第 \ref{ch:classic} 章中详细研究，作为更一般低维结构的先驱。在第 \ref{ch:autoencoding} 章中，我们将研究如何学习一般（例如非线性）低维分布的一致自编码。

\subsection{学习自洽表示}
请注意，在上述自编码目标中，需要评估解码数据 $\hat \X$ 与原始数据 $\X$ 的接近程度或一致性。这通常需要一些外部监督或知识来确定使用何种相似性度量。如果计算 $\hat \X$ 和 $\X$ 之间的相似性非常昂贵，甚至完全不可能或不可行，那么这将是一个问题。\footnote{假设一个人希望最小化两者之间的某种分布距离。} 请注意，在自然界中，动物能够在没有将估计值 $\hat \X$ 与数据空间中的真实值 $\X$ 进行比较的情况下，仅凭自身学习。它们甚至通常没有这种选择。

那么，一个系统如何在没有外部监督或比较的情况下自主学习呢？它们如何知道 $\hat \X$ 与 $\X$ 一致，即使没有直接比较它们？这导致了“闭合循环”的想法。事实证明，在第 \ref{ch:closed-loop} 章中我们将详细说明的温和条件下，为了确保 $\X$ 和 $\hat \X$ 一致，只需要将 $\hat \X$ 编码为 $\hat \Z$，并检查 $\Z$ 和 $\hat \Z$ 是否一致。我们将这种一致性称为 {\em 自洽性}，可以用以下图示说明：
\begin{equation}
    \X   \xrightarrow{\hspace{2mm} \mathcal{E}\hspace{2mm}} \Z  \xrightarrow{\hspace{2mm} \mathcal{D} \hspace{2mm}} \hat \X \xrightarrow{\hspace{2mm} \mathcal{E}\hspace{2mm}} \hat \Z,
    \label{eqn:closed-loop}
\end{equation}
我们将这个过程称为 {\em 闭合循环转录}，\footnote{受 DNA 和 RNA 或其他蛋白质之间的转录过程启发。} 如图 \ref{fig:closed-loop} 所示。

\begin{figure}[t]
\centering{\includegraphics[width=0.9\linewidth]{figures/diagrams_redu_gan_2.pdf}}
\caption{闭合循环转录的示意图。这里我们使用映射 $f$ 表示编码器 $\mathcal{E}$，使用 $g$ 表示解码器 $\mathcal{D}$。}  \label{fig:closed-loop}
\end{figure}

可以认为，任何自主智能体只需要学习观察数据 $\X$ 的自洽表示 $\Z$，因为检查原始数据空间（通常意味着外部世界）中的一致性要么太昂贵，要么甚至在物理上不可行。闭合循环公式允许通过仅依赖于内部（学习）特征 $\Z$ 的目标来学习最优编码和解码：
\begin{equation}
    \min d( \Z, \hat \Z) + \rho(\Z), 
   \label{eqn:closed-loop-objective}
\end{equation}
其中 $d( \Z, \hat \Z)$ 是相似性或距离度量，如我们将看到的，这可以更容易地计算。同样，$\rho(\Z)$ 是某种度量，促进 $\Z$ 的简约性和信息丰富性。令人惊讶的是，如我们在第 \ref{ch:closed-loop} 章中将看到的，在诸如 $\X$ 足够低维等相当温和的条件下，$(\Z, \hat \Z)$ 之间的自洽性意味着 $(\X, \hat \X)$ 之间的一致性！此外，我们还将看到，闭合循环系统将允许我们以 {\em 连续和增量} 的方式学习分布，\footnote{即一次学习一个类别，甚至一次学习一个样本。} 而不会遭受与开环模型相关的灾难性遗忘等问题。

\section{机器智能的理论与实践桥梁}
到目前为止，我们已经介绍了三种相关的框架，用于学习给定数据分布 $\X$ 的紧凑和结构化表示 $\Z$：
\begin{itemize}
\item 开放式编码 \eqref{eqn:open-loop};
\item 双向自编码 \eqref{eqn:auto-encoding};
\item 闭合循环转录 \eqref{eqn:closed-loop}。
\end{itemize}
在本书中，我们将系统地逐一研究这三个框架：
\begin{equation}
    \mbox{\textbf{开放式}} \; \Longrightarrow \; 
    \mbox{\textbf{双向}} \;  \Longrightarrow \; \mbox{\textbf{闭合循环}},
\end{equation}
分别在第 \ref{ch:representation}、\ref{ch:consistent} 和 \ref{ch:closed-loop} 章中。

在过去几年中，许多理论框架被提出和发展，以帮助理解深度网络。然而，许多理论未能提供与实际数据和任务上经验方法相匹配的可扩展解决方案。许多理论没有提供有用的指导，说明如何进一步改进实践。第 \ref{ch:applications} 章将展示本书中提出的框架如何有助于弥合理论与实践之间的差距。我们将提供令人信服的实验证据，证明从第一原理设计的网络和系统在图像分类、补全、分割和文本生成等任务上可以实现具有竞争力甚至更好的性能。

\paragraph{回到智能。}
正如我们在开头提到的，任何智能体的基本任务是从其感知的数据中学习可预测的信息。现在我们对这一任务的计算性质有了些许了解，应该意识到这是一个永无止境的过程，原因如下：
\begin{itemize}
    \item 从数据中学习到的当前知识，例如编码和解码方案，永远不可能真正最优。如果在预测新观察结果时仍然存在错误，总会有改进的空间。
    \item 到目前为止观察到的数据尚未涵盖所有可能的可预测信息。应该有机制来识别当前知识的不足，并必须学习和添加新的信息。
\end{itemize}

因此，智能不仅仅是提前收集所有数据并训练一个模型来记忆数据中的所有可预测信息。相反，它是关于配备能够不断改进当前信息并在需要时获取新信息的计算机制。也就是说，任何智能系统的根本特征之一是{\em 能够持续获取信息}（或知识）。我们可以将智能与信息（或知识）之间的关系表示如下：
\begin{equation}
\mbox{智能} = \frac{d\,}{dt} \mbox{信息}, \qquad 
\mbox{信息}  = \int_0^t \mbox{智能}\ dt.
\end{equation}
我们相信，上述闭环框架是一种通用机制，通过闭环反馈和游戏实现自我改进和自我学习。自然界中的所有智能体都在各个层面上利用闭环学习机制。其普遍性激发了早期尝试通过机器和计算机建模和模仿动物智能的研究，特别是诺伯特·维纳在 20 世纪 40 年代发起的控制论运动。

\end{document}